\documentclass[10pt,a4paper,reqno]{amsart}

\usepackage{amssymb,amsmath,amsthm,mathtools,bbm}
\usepackage[utf8]{inputenc}
\usepackage[T1]{fontenc}
\usepackage{lmodern}
\usepackage{enumitem}
\usepackage{hyperref}
\usepackage[capitalise]{cleveref}
\usepackage{natbib}
\usepackage{booktabs}
\usepackage{graphicx}
\usepackage{float}
\usepackage{siunitx}

% Path to figures produced by the Python notebooks
\graphicspath{{Code/notebooks/figures/}}

% Placeholder command: marks values to be updated from notebook output
\newcommand{\tbd}[1]{\textup{#1}}

\hypersetup{
  colorlinks=true,
  linkcolor=blue,
  citecolor=blue,
  urlcolor=blue,
  pdftitle={Pricing and Hedging of Renewable Energy Quanto Options in CARMA Models},
  pdfauthor={Vincent Gachon and Sven Karbach}
}
\setcitestyle{numbers,comma,open={[},close={]}}

\newtheorem{theorem}{Theorem}[section]
\newtheorem{proposition}[theorem]{Proposition}
\newtheorem{corollary}[theorem]{Corollary}
\newtheorem{lemma}[theorem]{Lemma}
\newtheorem{definition}[theorem]{Definition}
\newtheorem{remark}[theorem]{Remark}
\newtheorem{assumption}[theorem]{Assumption}

\newcommand{\R}{\mathbb{R}}
\newcommand{\C}{\mathbb{C}}
\newcommand{\E}{\mathbb{E}}
\newcommand{\Q}{\mathbb{Q}}
\newcommand{\PP}{\mathbb{P}}
\newcommand{\F}{\mathcal{F}}
\newcommand{\1}{\mathbbm{1}}
\newcommand{\dd}{\mathrm{d}}
\newcommand{\ii}{\mathrm{i}}
\newcommand{\ee}{\mathrm{e}}
\DeclareMathOperator{\Cov}{Cov}
\DeclareMathOperator{\Var}{Var}
\DeclareMathOperator{\diag}{diag}
\DeclareMathOperator{\Tr}{Tr}

\title[Pricing and Hedging of Quantos in CARMA Models]{Pricing and Hedging of Renewable Energy Quanto Options in CARMA Models}

\author{Vincent Gachon}
\address{Informatics Institute and Korteweg-de Vries Institute for Mathematics, University of Amsterdam, Science Park 105--107, 1098 XG Amsterdam, Netherlands}
\email{vincegachon@gmail.com}

\author{Sven Karbach}
\address{Informatics Institute and Korteweg-de Vries Institute for Mathematics, University of Amsterdam, Science Park 105--107, 1098 XG Amsterdam, Netherlands}
\email{sven@karbach.org}

\keywords{quanto option, CARMA process, multivariate CARMA, incomplete market, variance-optimal hedging, energy derivatives}

\begin{document}

\begin{abstract}
We study quanto claims whose payoff depends on a tradable energy factor and on a
second factor that captures volumetric, weather, foreign-exchange, or
cross-market risk. The paper develops a coupled continuous-time autoregressive
moving average (CARMA) framework in which the log-price factor and the secondary
factor are driven by Lévy noises linked through a directed coupling channel.
This specification preserves the interpretable low-dimensional structure of
univariate CARMA models while retaining a tractable joint characteristic
function and an explicit state-space representation. We show that the model
embeds naturally into a multivariate CARMA formulation, which clarifies when
the coupled specification is sufficient and when a full multivariate treatment
is preferable.

On the pricing side we formulate the problem on a filtered space with domestic
numeraire and an exogenously chosen pricing measure, which is essential because
the weather or quantity factor is typically non-tradable. We derive conditional
transform formulas, forward-price representations, and a Fourier valuation
result for a broad class of European quanto payoffs. For the practically
important case of weak Brownian coupling, we obtain a first-order quanto
adjustment that isolates the contribution of the dependence channel to option
values. On the hedging side we emphasize market incompleteness and characterize
the variance-optimal hedge through the Galtchouk--Kunita--Watanabe (GKW)
decomposition. This yields a transparent separation between the hedgeable energy
price exposure and the residual hybrid risk that cannot be eliminated without
additional tradable instruments.
\end{abstract}

\maketitle

\section{Introduction}\label{sec:introduction}

Quanto-type contracts in energy and climate-sensitive markets are hybrid claims:
their value depends simultaneously on a tradable price signal and on a second
factor that captures volume, temperature, irradiance, wind availability,
congestion, or settlement-currency effects. In retail gas, power, and renewable
generation portfolios this second factor is often the economically dominant
source of non-standard exposure because demand and production volumes react to
weather conditions while revenues remain tied to commodity prices. Energy quanto
contracts are therefore designed to hedge joint price and quantity risk rather
than either component in isolation; see \citet{Ben15} for the classical
futures-based approach and \citet{Alf24} for a recent coupled
electricity--temperature model.

The modeling challenge is structural. The secondary factor is often not
tradable, so no-arbitrage arguments alone do not determine a unique pricing
measure. At the same time, the dependence between the price factor and the
secondary factor matters materially for both valuation and hedging. Simple
bivariate Ornstein--Uhlenbeck specifications are appealing because they lead to
closed formulas, but they impose overly rigid autocovariance structures and do
not exploit the richer memory patterns observed in power prices, renewable
generation, temperature, and other environmental series. In electricity markets,
Lévy-driven CARMA (continuous-time autoregressive moving average) dynamics are
now a standard way of capturing mean reversion, serial dependence, and spikes
while maintaining tractable state-space representations; see \citet{Bro01},
\citet{Ben13}, and \citet{Ben14}. A CARMA$(p,q)$ process is the
continuous-time analogue of an ARMA$(p,q)$ time series: the autoregressive
polynomial of degree $p$ governs the mean-reversion speed and spectral shape,
while the moving-average polynomial of degree $q < p$ shapes the impulse
response. Driving the state equation with a Lévy process allows heavy tails,
jumps, and skewness without sacrificing the linear state-space structure. The
same state-space machinery is attractive in hybrid derivative problems because
it supports transform methods, filtering, and dynamic hedging computations.

The present paper develops a pricing and hedging framework for energy quanto
options driven by coupled CARMA factors. The word ``quanto'' is used here in
the broad sense common in energy markets: the payoff depends on two underlying
indices, one typically tradable and one often non-tradable. The classical
foreign-exchange quanto is recovered when the secondary factor is itself a
tradable exchange rate. We formulate the problem in a way that covers both
interpretations. The starting point is a coupled state-space specification in
which the secondary factor drives its own CARMA dynamics and enters the price
factor through a directional coupling term. This construction preserves
parsimony and economic interpretation: the secondary factor can represent a
causal weather, load, or FX channel affecting the price process without forcing
the entire system into a fully symmetric multivariate model.

The contribution of the paper is fivefold.
\begin{enumerate}[label=(\roman*)]
  \item We formulate a coupled Lévy-driven CARMA model for the joint dynamics of
    a tradable energy factor and a secondary risk factor, and we derive its
    state-space, cross-covariance, and characteristic-function representations.
  \item We show that the coupled specification is a structured special case of a
    multivariate CARMA model in the sense of \citet{Mar07}; this clarifies when
    the lower-dimensional coupled form is appropriate and when a full MCARMA
    formulation is preferable. In particular, the multivariate viewpoint becomes
    compelling once several tradeable commodities, locations, or currencies are
    modeled jointly, or when feedback effects are empirically relevant.
  \item We develop a valuation framework under a domestic numeraire and a chosen
    pricing measure. For a broad class of European quanto payoffs we derive a
    Fourier representation in terms of the conditional joint characteristic
    function. We also obtain a first-order quanto adjustment for weak coupling,
    which isolates the marginal pricing impact of the dependence channel.
  \item We analyze hedging in the incomplete market generated by a non-tradable
    secondary factor. The appropriate benchmark is not perfect replication but
    variance-optimal hedging. We characterize the hedge by the
    Galtchouk--Kunita--Watanabe decomposition and show how the residual risk
    measures the irreducible hybrid exposure.
  \item We calibrate the coupled CARMA framework to hourly German day-ahead
    electricity prices and temperature observations (2020--2025), fit
    Normal Inverse Gaussian driving noise, estimate the coupling parameter from
    cross-covariance data, validate the Fourier pricing formula against Monte
    Carlo simulation, and compare the resulting variance-optimal hedge against
    an OU benchmark in an out-of-sample 2025 backtest.
\end{enumerate}

Relative to the literature on energy quantos, the novelty is therefore not a
single new numerical scheme in isolation, but a coherent mathematical framework
that links CARMA dynamics, measure selection, transform pricing, and
incomplete-market hedging in a single model class, together with an empirical
demonstration on real hourly data.

The remainder of the paper is organized as follows.
\Cref{sec:carma-preliminaries} recalls the state-space representation of
Lévy-driven CARMA factors. \Cref{sec:coupled-carma-models} introduces the
coupled model, derives its joint transform, and relates it to multivariate
CARMA. \Cref{sec:pricing-quanto-options} develops the pricing framework,
including Fourier valuation and a weak-coupling expansion. Hedging is discussed
in \Cref{sec:hedging-strategies}. \Cref{sec:calibration} presents the empirical
calibration to German electricity and temperature data. \Cref{sec:numerical}
reports the pricing and hedging numerical experiments. \Cref{sec:limitations}
highlights limitations and extensions, and \Cref{sec:conclusion} concludes.

\section{CARMA Preliminaries}\label{sec:carma-preliminaries}

We work on a filtered probability space
$(\Omega,\F,(\F_t)_{t \ge 0},\mathbb{M})$ satisfying the usual conditions,
where $\mathbb{M}$ stands for either the historical measure $\PP$ or a
pricing measure $\Q$. All processes below are assumed càdlàg and adapted.

\begin{definition}\label{def:carma-factor}
Let $L=(L_t)_{t \ge 0}$ be an $m$-dimensional Lévy process under
$\mathbb{M}$, let $A \in \R^{p \times p}$, $B \in \R^{p \times m}$, and
let $c \in \R^p$. We call the output process
$$
Y_t = c^\top Z_t,\qquad t \ge 0,
$$
a \emph{Lévy-driven CARMA factor in state-space form} if the state process
$(Z_t)_{t \ge 0}$ solves
\begin{align}\label{eq:state_general}
  \dd Z_t = A Z_t\,\dd t + B\,\dd L_t.
\end{align}
When $A$ is chosen as the companion matrix associated with a stable
autoregressive polynomial of degree $p$, and $B$, $c$ are determined by
the moving-average polynomial of degree $q < p$,
\eqref{eq:state_general} is the standard state-space
representation of a CARMA$(p,q)$ process; see \citet{Bro01}. Since the driving
noise $L$ is $m$-dimensional while the output $Y_t$ is scalar, this is a
\emph{multi-input single-output} (MISO) state-space system; the standard
CARMA$(p,q)$ process is recovered as the special case $m=1$ with companion
matrix $A$. In the present paper the state-space form is the primary object
because it is the natural representation for pricing and hedging. The scalar
output $Y_t = c^\top Z_t$ has the rational transfer function
$G(z) = c^\top(zI-A)^{-1}B$, whose numerator and denominator polynomials
correspond respectively to the moving-average and autoregressive components.
\end{definition}

\begin{assumption}\label{ass:stableA}
All eigenvalues of $A$ have strictly negative real part.
\end{assumption}

\begin{proposition}[Variation of constants and stationarity]
\label{prop:variation_constants}
Under \Cref{ass:stableA}, the state equation \eqref{eq:state_general} has the
variation-of-constants representation
\begin{align}\label{eq:variation_constants}
  Z_t = \ee^{A(t-s)} Z_s + \int_s^t \ee^{A(t-u)} B\,\dd L_u,
  \qquad 0 \le s \le t.
\end{align}
If $\mathbb{E}^{\mathbb{M}}[\log(1+\|L_1\|)] < \infty$, then there exists a
unique strictly stationary solution given by
\begin{align}\label{eq:stationary_solution}
  Z_t = \int_{-\infty}^t \ee^{A(t-u)} B\,\dd L_u,
  \qquad t \ge 0.
\end{align}
\end{proposition}

\begin{proof}
Equation \eqref{eq:variation_constants} is the standard linear SDE solution.
Under \Cref{ass:stableA}, $\|\ee^{At}\|$ decays exponentially. The condition
$\mathbb{E}^{\mathbb{M}}[\log(1+\|L_1\|)] < \infty$ is the minimal
integrability requirement for the improper stochastic integral in
\eqref{eq:stationary_solution} to converge almost surely and in $L^1$;
see \citet{Bro01}. Under these conditions the integral gives the unique strictly
stationary solution of \eqref{eq:state_general}.
\end{proof}

For transform pricing we need the conditional characteristic function of the
state-space output.

\begin{proposition}[Conditional characteristic function]
\label{prop:cf_single_factor}
Let $\psi_{\mathbb{M}}:\R^m \to \C$ denote the characteristic exponent of
$L$ under $\mathbb{M}$ (also called the Lévy symbol or cumulant generating
function on the imaginary line), so that
$$
\E^{\mathbb{M}}\!\left[\ee^{\ii\langle u,L_t\rangle}\right]
= \ee^{t\psi_{\mathbb{M}}(u)},\qquad u \in \R^m;
$$
see \citet{Con04} for background on Lévy processes and their characteristic
exponents.
Then for $0 \le t \le T$ and $u \in \R$,
\begin{align}\label{eq:single_factor_cf}
  \E^{\mathbb{M}}\!\left[\ee^{\ii u Y_T}\mid \F_t\right]
  = \exp\!\left(
      \ii u\, c^\top \ee^{A(T-t)} Z_t
      + \int_t^T \psi_{\mathbb{M}}\!\left(
          u\, B^\top \ee^{A^\top(T-s)} c
        \right)\dd s
    \right).
\end{align}
\end{proposition}

\begin{proof}
By \eqref{eq:variation_constants},
$$
Y_T
= c^\top \ee^{A(T-t)} Z_t
  + \int_t^T c^\top \ee^{A(T-s)} B\,\dd L_s.
$$
The stochastic integral over $[t,T]$ is independent of $\F_t$, and its
characteristic function is determined by the Lévy exponent. This yields
\eqref{eq:single_factor_cf}.
\end{proof}

\begin{remark}[Two-sided Lévy extension]
\label{rem:two_sided_levy}
The improper stochastic integral in \eqref{eq:stationary_solution} requires $L$
to be defined on the whole line $\mathbb{R}$. A canonical construction is to take
an independent copy $L^{(2)}$ of $L$ and set
$$
L_t
= \begin{cases}
    L_t^{(1)}, & t \ge 0,\\
    -L_{(-t)-}^{(2)}, & t < 0,
  \end{cases}
$$
where $L^{(1)}$ is the original one-sided process and $L_{t-}^{(2)} =
\lim_{s\nearrow t}L_s^{(2)}$ denotes the left limit.
Under this convention the integral in \eqref{eq:stationary_solution} is a
well-defined $L^1$-limit whenever $\mathbb{E}[\log(1+\|L_1\|)]<\infty$ and
$A$ is stable; see \citet{Bro01} for details.
\end{remark}

\begin{remark}
The state-space form is not just a notational convenience. It makes explicit the
transfer kernel $s \mapsto c^\top \ee^{As} B$, which is the object needed for
transform methods, Monte Carlo simulation, and the derivation of covariance
equations. It also clarifies why CARMA models remain analytically tractable
under Lévy forcing: linearity reduces pricing questions to matrix exponentials
and Lévy cumulants.
\end{remark}

\section{Coupled CARMA Models for Quanto Risk}\label{sec:coupled-carma-models}

\noindent\textit{Convention.}
From this section onward, all processes, expectations, and characteristic
functions are understood to be under the pricing measure $\Q$ unless
explicitly indicated otherwise.  The superscript $\mathbb{M}$ is retained
only in statements that hold simultaneously under $\PP$ and $\Q$, as in the
second-order structure of \Cref{sec:coupled-carma-models}.
State vectors such as $Z_t^{X,\Q}$ carry the measure superscript in the
model equations where the distinction matters; in the pricing and hedging
sections the superscript is dropped for readability since $\Q$ is
unambiguous.

\subsection{Model specification}

We distinguish a \emph{tradable} energy factor from a \emph{secondary} factor.
The secondary factor may be temperature, wind, solar irradiance, demand, an FX
rate, or any other quantity entering the payoff. Denote by $X_t$ the
deseasonalized log-price factor and by $Y_t$ the deseasonalized secondary
factor. The observed spot price is $S_t = \ee^{\Lambda_S(t)+X_t}$, where
$\Lambda_S$ is a deterministic seasonal component, and the observed secondary
index is $\Lambda_Y(t)+Y_t$, with deterministic component $\Lambda_Y$.

The dynamics under a generic measure $\mathbb{M} \in \{\PP,\Q\}$ are
\begin{align}\label{eq:coupled_model}
\begin{aligned}
  \dd Z_t^{Y,\mathbb{M}}
    &= A_Y^{\mathbb{M}} Z_t^{Y,\mathbb{M}}\,\dd t
       + B_Y\,\dd L_t^{Y,\mathbb{M}},\\
  \dd Z_t^{X,\mathbb{M}}
    &= A_X^{\mathbb{M}} Z_t^{X,\mathbb{M}}\,\dd t
       + B_X\,\dd L_t^{X,\mathbb{M}}
       + \Gamma\,\dd L_t^{Y,\mathbb{M}},
\end{aligned}
\end{align}
with output equations
\begin{align}\label{eq:outputs_coupled}
  X_t = c_X^\top Z_t^{X,\mathbb{M}},
  \qquad
  Y_t = c_Y^\top Z_t^{Y,\mathbb{M}}.
\end{align}
Here the state vectors satisfy $Z^{X,\mathbb{M}} \in \R^{p_X}$ and
$Z^{Y,\mathbb{M}} \in \R^{p_Y}$, and the Lévy drivers $L^{X,\mathbb{M}} \in
\R^{m_X}$ and $L^{Y,\mathbb{M}} \in \R^{m_Y}$ are independent.
The \emph{drift matrices} $A_X^{\mathbb{M}} \in \R^{p_X \times p_X}$ and
$A_Y^{\mathbb{M}} \in \R^{p_Y \times p_Y}$ are measure-dependent (they absorb
the effect of Girsanov or Esscher changes). The \emph{noise-loading matrices}
$B_X \in \R^{p_X \times m_X}$, $B_Y \in \R^{p_Y \times m_Y}$, the
\emph{coupling matrix} $\Gamma \in \R^{p_X \times m_Y}$, and the
\emph{output vectors} $c_X \in \R^{p_X}$, $c_Y \in \R^{p_Y}$ are
structural parameters that do not depend on the choice of measure. This is
important: the observation equation defines the \emph{same} physical
processes $X_t$ and $Y_t$ under every equivalent measure.

\begin{assumption}\label{ass:coupled_model}
Throughout the paper we impose the following conditions.
\begin{enumerate}[label=(\alph*)]
  \item\label{ass:stability} The drift matrices $A_X^{\mathbb{M}}$ and
    $A_Y^{\mathbb{M}}$ are stable (all eigenvalues have strictly negative real
    part) under every measure $\mathbb{M}$ considered.
  \item\label{ass:moments} The Lévy drivers $L^{X,\mathbb{M}}$ and
    $L^{Y,\mathbb{M}}$ are independent and admit exponential moments: there
    exist $\theta, \rho > 0$ such that
    $$
      \mathbb{E}^{\mathbb{M}}\!\left[
        \ee^{\,\theta\|L_1^{X,\mathbb{M}}\| + \rho\|L_1^{Y,\mathbb{M}}\|}
      \right] < \infty,
    $$
    which guarantees that the cumulant exponents are finite on a non-empty open
    strip $\mathcal{S} \subset \R^{m_X} \times \R^{m_Y}$ around the origin,
    enabling forward-price computations and Fourier inversion.
  \item\label{ass:pricing_measure_main} There exists an equivalent probability
    measure $\Q$ such that:
    \begin{enumerate}[label=(\roman*)]
      \item under $\Q$, the factor dynamics retain the coupled state-space form
        \eqref{eq:coupled_model} with the same noise-loading matrices $B_X,
        B_Y, \Gamma$ and output vectors $c_X, c_Y$, but possibly different
        drift matrices $A_X^\Q, A_Y^\Q$ and Lévy triplets;
      \item the drivers $L^{X,\Q}$ and $L^{Y,\Q}$ remain independent under
        $\Q$;
      \item discounted tradable forward prices formed from the energy factor are
        $\Q$-martingales.
    \end{enumerate}
    A standard mechanism satisfying (i)--(ii) is the Esscher transform applied
    to the joint Lévy driver $(L^X,L^Y)$ (see \Cref{rem:esscher_preservation}
    below). Condition~(iii) imposes no-arbitrage consistency and determines the
    drift adjustment of the energy factor under~$\Q$.
\end{enumerate}
\end{assumption}

\begin{remark}\label{rem:esscher_preservation}
The Esscher transform preserves the coupled structure. If $L = (L^X, L^Y)$
has independent components and one applies an exponential tilt with parameter
vector $(\vartheta_X, \vartheta_Y) \in \mathcal{S}$, the resulting measure
$\Q^{\vartheta}$ modifies the Lévy triplet of each component independently.
In particular, $L^{X,\Q^\vartheta}$ and $L^{Y,\Q^\vartheta}$ remain
independent, and the linear state-space dynamics \eqref{eq:coupled_model} are
preserved with the same structural matrices $B_X, B_Y, \Gamma$. The drift
matrices $A_X, A_Y$ change only if the Esscher tilt interacts with the
compensator of a jump part (which adds a drift correction); in the Gaussian
case, the Esscher tilt reduces to the classical Girsanov drift shift. See
\citet{Ben14} for details in the energy CARMA context.

\paragraph{Drift adjustment and martingale calibration.}
Under the Esscher measure $\Q^\vartheta$ the cumulant exponent of $L^X$
shifts to
$$
  \kappa_X^{\Q}(u)
  = \kappa_X^{\PP}(u + \vartheta_X) - \kappa_X^{\PP}(\vartheta_X),
  \qquad u \in \mathcal{S}_X - \vartheta_X.
$$
The tilt modifies the mean of $L^{X,\Q}$, which in the state equation
\eqref{eq:coupled_model} produces a drift correction that enters the last
row of the companion matrix: in the CARMA$(p_X,q_X)$ parametrisation one
has
$$
  A_X^{\Q}
  = A_X^{\PP}
    + B_X
      \bigl(\kappa_X^{\PP}{}'\!(\vartheta_X)
            - \kappa_X^{\PP}{}'\!(0)\bigr)
      c_X^\top,
$$
where $\kappa_X^{\PP}{}'$ denotes the gradient of the cumulant exponent.

The no-arbitrage condition~\Cref{ass:coupled_model}\ref{ass:pricing_measure_main}(iii)
requires the discounted spot $\ee^{-r_t^d t}S_t$ to be a $\Q$-local martingale,
where $S_t = \ee^{\Lambda_S(t)+X_t}$.  Using the forward price representation
\eqref{eq:forward_formula} this yields one scalar equation in $\vartheta_X$:
\begin{align}\label{eq:esscher_calibration}
  \kappa_X^{\Q}\!\bigl(c_X^\top(-A_X^{\Q})^{-1}B_X\bigr)
  + \kappa_Y^{\Q}\!\bigl(c_X^\top(-A_X^{\Q})^{-1}\Gamma\bigr)
  = r^d - \dot\Lambda_S,
\end{align}
where $r^d$ is the (constant) domestic short rate from
\Cref{ass:interest_rate} and $\dot\Lambda_S = \dd\Lambda_S/\dd t$ is the
seasonal trend drift.
For NIG drivers both $\kappa_X^{\Q}$ and $\kappa_Y^{\Q}$ are analytic and
strictly monotone in their respective tilt parameters, so
\eqref{eq:esscher_calibration} has a unique solution $\vartheta_X$ whenever
the right-hand side lies in the interior of the range of the left-hand side.
The tilt parameter $\vartheta_Y$ for the secondary (non-tradable) factor
remains a free parameter that encodes the market price of weather or quantity
risk; in the absence of calibration instruments for that factor one may set
$\vartheta_Y = 0$, which corresponds to pricing under the historical law of
the secondary factor.
\end{remark}

\begin{remark}
The directional structure in \eqref{eq:coupled_model} is deliberate. It is
appropriate when one wants the weather or quantity factor to affect the price
factor but not conversely at the same horizon. This is natural for many
renewable applications: temperature, wind, or irradiation shocks move expected
production or demand and thereby affect power prices. The model remains
parsimonious because dependence is introduced through the single matrix
$\Gamma$ instead of a fully unrestricted multivariate dynamics.
\end{remark}

\begin{remark}
The classical foreign-exchange quanto fits into the same notation. If the
secondary factor is the log FX rate and is itself tradable, then the pair
$(X,Y)$ describes joint commodity and currency risk. The main difference is
conceptual rather than mathematical: in the tradable FX case the pricing measure
is tied to domestic and foreign numeraires, whereas in the weather case the
secondary factor is typically non-tradable and the market is incomplete.
\end{remark}

\subsection{Second-order structure}

When the Lévy drivers are square-integrable, the coupled model admits an
explicit stationary cross-covariance equation. This is useful both for
calibration and for understanding the sign and decay of the quanto effect.

\begin{proposition}[Stationary cross-covariance]
\label{prop:cross_covariance}
Assume that under $\mathbb{M}$ the processes $L^{X,\mathbb{M}}$ and
$L^{Y,\mathbb{M}}$ are independent, square-integrable, and centered, with
$$
\Sigma_Y^{\mathbb{M}} = \Var^{\mathbb{M}}(L_1^{Y,\mathbb{M}}).
$$
Let $(Z_t^{X,\mathbb{M}},Z_t^{Y,\mathbb{M}})$ denote the stationary solution
of \eqref{eq:coupled_model}. Then the matrix
$$
M^{\mathbb{M}} = \Cov^{\mathbb{M}}(Z_t^{X,\mathbb{M}}, Z_t^{Y,\mathbb{M}})
$$
is time-independent and solves the Sylvester equation
\begin{align}\label{eq:sylvester_cross_cov}
  A_X^{\mathbb{M}} M^{\mathbb{M}}
  + M^{\mathbb{M}} (A_Y^{\mathbb{M}})^\top
  + \Gamma \Sigma_Y^{\mathbb{M}} B_Y^\top
  = 0.
\end{align}
For $h \ge 0$,
\begin{align}\label{eq:lagged_cross_cov}
  \Cov^{\mathbb{M}}(X_{t+h},Y_t)
  = c_X^\top
      \ee^{A_X^{\mathbb{M}} h}
      M^{\mathbb{M}}
      c_Y.
\end{align}
\end{proposition}

\begin{proof}
Let $N_t \coloneqq Z_t^{X,\mathbb{M}}(Z_t^{Y,\mathbb{M}})^\top$. By Itô's
formula for semimartingales,
\begin{align*}
  \dd N_t
  &= (\dd Z_t^{X,\mathbb{M}})(Z_t^{Y,\mathbb{M}})^\top
     + Z_t^{X,\mathbb{M}}(\dd Z_t^{Y,\mathbb{M}})^\top
     + \dd [Z^{X,\mathbb{M}}, Z^{Y,\mathbb{M}}]_t^c + \Delta N_t,
\end{align*}
where $[{\cdot},{\cdot}]^c$ denotes the continuous quadratic covariation.
The continuous covariation between the $X$-state and the $Y$-state receives
a contribution from the shared noise $L^{Y,\mathbb{M}}$ entering both
equations through $B_Y$ and $\Gamma$, and a zero contribution from
$L^{X,\mathbb{M}}$ (which is independent of $L^{Y,\mathbb{M}}$).
For the jump covariation: since the only shared noise between $Z^{X,\mathbb{M}}$
and $Z^{Y,\mathbb{M}}$ is $L^{Y,\mathbb{M}}$, simultaneous jumps at time $s$
satisfy $\Delta Z_s^{X,\mathbb{M}} = \Gamma\,\Delta L_s^{Y,\mathbb{M}}$ and
$\Delta Z_s^{Y,\mathbb{M}} = B_Y\,\Delta L_s^{Y,\mathbb{M}}$.
Therefore the expected rate of the jump covariation is
$$
  \mathbb{E}^{\mathbb{M}}\!\left[
    \frac{\dd}{\dd t}\sum_{s \le t}
      \Delta Z_s^{X,\mathbb{M}}\bigl(\Delta Z_s^{Y,\mathbb{M}}\bigr)^\top
  \right]
  = \Gamma\,
    \mathbb{E}^{\mathbb{M}}\!\bigl[
      \Delta L^{Y,\mathbb{M}}\bigl(\Delta L^{Y,\mathbb{M}}\bigr)^\top
    \bigr]\,
    B_Y^\top
  = \Gamma\,\Sigma_Y^{\mathbb{M}}\,B_Y^\top,
$$
where
$\mathbb{E}^{\mathbb{M}}\bigl[\sum_{s\le 1}\Delta L^Y_s(\Delta L^Y_s)^\top\bigr]
= \int_{\mathbb{R}^{m_Y}} x x^\top\,\nu_Y(dx) = \Sigma_Y^{\mathbb{M}}$
by the Lévy--Itô decomposition and the square-integrability in
\Cref{ass:coupled_model}\ref{ass:moments}.
Substituting the state equations \eqref{eq:coupled_model} and taking
expectations in the stationary regime
($\mathbb{E}[\dd N_t/\dd t] = 0$) gives
$$
  A_X^{\mathbb{M}} M^{\mathbb{M}}
  + M^{\mathbb{M}} (A_Y^{\mathbb{M}})^\top
  + \Gamma \Sigma_Y^{\mathbb{M}} B_Y^\top
  = 0,
$$
which is \eqref{eq:sylvester_cross_cov}. Equivalently, the solution can be
expressed as
$M^{\mathbb{M}} = \int_0^\infty \ee^{A_X^{\mathbb{M}} \tau}\,
\Gamma \Sigma_Y^{\mathbb{M}} B_Y^\top\,
\ee^{(A_Y^{\mathbb{M}})^\top \tau}\,\dd\tau$,
which is finite under \Cref{ass:coupled_model}\ref{ass:stability}. The lagged
formula \eqref{eq:lagged_cross_cov} follows from the Markov property and the
variation-of-constants formula \eqref{eq:variation_constants}.
\end{proof}

\begin{remark}[Unique solvability of the Sylvester equation]
\label{rem:sylvester_unique}
The equation $A_X^{\mathbb{M}} M + M(A_Y^{\mathbb{M}})^\top = C$ has a unique
solution $M \in \mathbb{R}^{p_X \times p_Y}$ for every right-hand side $C$ if
and only if $A_X^{\mathbb{M}}$ and $-(A_Y^{\mathbb{M}})^\top$ share no common
eigenvalue, i.e.\ $\lambda_i(A_X^{\mathbb{M}}) + \lambda_j(A_Y^{\mathbb{M}}) \ne
0$ for all $i,j$.
Under \Cref{ass:coupled_model}\ref{ass:stability} all eigenvalues of
$A_X^{\mathbb{M}}$ and $A_Y^{\mathbb{M}}$ have strictly negative real parts, so
$\mathrm{Re}\bigl(\lambda_i(A_X^{\mathbb{M}}) + \lambda_j(A_Y^{\mathbb{M}})\bigr)
< 0$ for every pair $(i,j)$; in particular the sum is never zero.
Hence \eqref{eq:sylvester_cross_cov} has a unique solution, given explicitly by
$$
  M^{\mathbb{M}}
  = \int_0^\infty
      \ee^{A_X^{\mathbb{M}}\tau}\,
      \Gamma\Sigma_Y^{\mathbb{M}} B_Y^\top\,
      \ee^{(A_Y^{\mathbb{M}})^\top\tau}
    \,\dd\tau,
$$
which is finite because both matrix exponentials decay to zero exponentially fast
under \Cref{ass:coupled_model}\ref{ass:stability}.
\end{remark}

\begin{remark}
The cross-covariance is governed by a linear matrix equation. This is a genuine
advantage of coupled CARMA models over ad hoc bivariate specifications: the sign
and decay of dependence are encoded in the loading matrix $\Gamma$, the noise
variance of the secondary factor, and the mean-reversion matrices
$A_X, A_Y$. In particular, one can separate the contribution of persistence
from the contribution of instantaneous coupling.
\end{remark}

\begin{remark}[Spectral viewpoint]
The transfer functions associated with \eqref{eq:coupled_model} are
\begin{align*}
  G_X(z) &= (c_X)^\top (zI-A_X)^{-1} B_X,\\
  G_{XY}(z) &= (c_X)^\top (zI-A_X)^{-1} \Gamma,\\
  G_Y(z) &= (c_Y)^\top (zI-A_Y)^{-1} B_Y.
\end{align*}
Hence the coupled model retains the rational spectral structure familiar from
univariate CARMA models. This is useful for order selection and for matching
empirical autocovariances or periodograms in calibration.
\end{remark}

\subsection{Conditional transform and forward prices}

We now derive the key pricing object: the conditional joint transform of
$(X_T,Y_T)$, i.e.\ the deseasonalized factor pair.

For $0 \le t \le T$, define
\begin{align*}
  \overline{X}_t(T)
    &\coloneqq c_X^\top \ee^{A_X^\Q (T-t)} Z_t^{X,\Q},\\
  \overline{Y}_t(T)
    &\coloneqq c_Y^\top \ee^{A_Y^\Q (T-t)} Z_t^{Y,\Q},
\end{align*}
as well as the deterministic loading kernels (which depend only on $T-s$,
not on $t$ separately)
\begin{align}
  \beta_X(s;T)
    &\coloneqq B_X^\top \ee^{(A_X^\Q)^\top (T-s)} c_X,
    \label{eq:betaX}\\
  \gamma_X(s;T)
    &\coloneqq \Gamma^\top \ee^{(A_X^\Q)^\top (T-s)} c_X,
    \label{eq:gammaX}\\
  \beta_Y(s;T)
    &\coloneqq B_Y^\top \ee^{(A_Y^\Q)^\top (T-s)} c_Y.
    \label{eq:betaY}
\end{align}
Let $\psi_X^\Q$ and $\psi_Y^\Q$ denote the characteristic exponents of
$L^{X,\Q}$ and $L^{Y,\Q}$.

\begin{proposition}[Joint transform]
\label{prop:joint_cf}
Under \Cref{ass:coupled_model}, for any $u,v \in \R$ and $0 \le t \le T$,
\begin{align}
  \varphi_t(u,v;T)
  &\coloneqq \E^\Q\!\left[
      \ee^{\ii u X_T + \ii v Y_T}
      \mid \F_t
    \right] \nonumber\\
  &= \exp\Bigg(
      \ii u\, \overline{X}_t(T)
      + \ii v\, \overline{Y}_t(T)
      + \int_t^T \psi_X^\Q\!\big(u\,\beta_X(s;T)\big)\,\dd s
      \nonumber\\
      &\hspace{3.6cm}
      + \int_t^T \psi_Y^\Q\!\big(
          u\,\gamma_X(s;T)
          + v\,\beta_Y(s;T)
        \big)\,\dd s
    \Bigg).
    \label{eq:joint_cf}
\end{align}
\end{proposition}

\begin{proof}
By variation of constants,
\begin{align*}
  X_T
  &= \overline{X}_t(T)
     + \int_t^T c_X^\top \ee^{A_X^\Q(T-s)} B_X\,\dd L_s^{X,\Q}
     + \int_t^T c_X^\top \ee^{A_X^\Q(T-s)} \Gamma\,\dd L_s^{Y,\Q},\\
  Y_T
  &= \overline{Y}_t(T)
     + \int_t^T c_Y^\top \ee^{A_Y^\Q(T-s)} B_Y\,\dd L_s^{Y,\Q}.
\end{align*}
The increment of $L^{X,\Q}$ over $[t,T]$ is independent of both $\F_t$
and $L^{Y,\Q}$, so the conditional characteristic function factorizes into two
Lévy transform terms. Evaluating those terms with the corresponding
characteristic exponents yields \eqref{eq:joint_cf}.
\end{proof}

The joint transform immediately produces the forward price of the energy factor.
Let $\kappa_X^\Q$ and $\kappa_Y^\Q$ denote the cumulant exponents of
$L^{X,\Q}$ and $L^{Y,\Q}$, defined on the set of arguments for which
exponential moments exist.

\begin{corollary}[Forward price representation]
\label{cor:forward_price}
Let $\mathcal{S}_X$ and $\mathcal{S}_Y$ denote the strips of regularity of
the cumulant exponents $\kappa_X^\Q$ and $\kappa_Y^\Q$, i.e.\ the sets of
real arguments on which the exponential moments are finite. Assume that for
all $s \in [t,T]$, $\beta_X(s;T) \in \operatorname{int}(\mathcal{S}_X)$
and $\gamma_X(s;T) \in \operatorname{int}(\mathcal{S}_Y)$. A sufficient
condition is that the kernel bounds
$\sup_{s \in [t,T]}\|\beta_X(s;T)\|$ and
$\sup_{s \in [t,T]}\|\gamma_X(s;T)\|$ are smaller than the radii of
$\mathcal{S}_X$ and $\mathcal{S}_Y$; the kernel bounds are finite by
stability of $A_X^\Q$ and $A_Y^\Q$, while the strip radii are determined
by the Lévy measures.
Then the forward price for physical delivery at time $T$ is
\begin{align}
  F(t,T)
  &= \E^\Q[S_T \mid \F_t] \nonumber\\
  &= \exp\Bigg(
      \Lambda_S(T)
      + \overline{X}_t(T)
      + \int_t^T \kappa_X^\Q\!\big(\beta_X(s;T)\big)\,\dd s
      + \int_t^T \kappa_Y^\Q\!\big(\gamma_X(s;T)\big)\,\dd s
    \Bigg).
    \label{eq:forward_formula}
\end{align}
\end{corollary}

\begin{proof}
Because $S_T = \ee^{\Lambda_S(T)+X_T}$, the forward price is
$\E^\Q[S_T\mid\F_t] = \ee^{\Lambda_S(T)} \E^\Q[\ee^{X_T}\mid\F_t]$.
Setting $u = -\ii$ and $v = 0$ in \eqref{eq:joint_cf} requires that
$\varphi_t(-\ii,0;T)$ is well-defined, i.e.\ that the argument $-\ii$ lies in the
domain of analyticity of $u \mapsto \varphi_t(u,0;T)$.
By \eqref{eq:joint_cf} this reduces to requiring
$\beta_X(s;T) \in \mathrm{int}(\mathcal{S}_X)$ and
$\gamma_X(s;T) \in \mathrm{int}(\mathcal{S}_Y)$ for all $s \in [t,T]$,
where we use the cumulant-generating function via $\psi(-\ii\,w) = \kappa(w)$ for
real $w$ in the strip.
We verify these inclusions as follows.
Under \Cref{ass:coupled_model}\ref{ass:stability} the kernels satisfy
$$
  \|\beta_X(s;T)\|
  \le \|B_X\|\,\|c_X\|\,
      \ee^{\mathrm{Re}(\lambda_{\max}(A_X^\Q))(T-s)},
  \quad
  \|\gamma_X(s;T)\|
  \le \|\Gamma\|\,\|c_X\|\,
      \ee^{\mathrm{Re}(\lambda_{\max}(A_X^\Q))(T-s)},
$$
so both are uniformly bounded on $[t,T]$ for any finite $T - t$.
Under \Cref{ass:coupled_model}\ref{ass:moments} the cumulant strips
$\mathcal{S}_X$ and $\mathcal{S}_Y$ each contain a ball of strictly positive
radius $r_X, r_Y > 0$ around the origin
(since $\int_{\|x\|>1}\ee^{\theta\|x\|}\nu_X(dx) < \infty$ for
$\theta \in (-\theta_0,\theta_0)$ with $\theta_0 \ge \min(\theta,\rho)$
from \Cref{ass:coupled_model}\ref{ass:moments}).
Sufficient conditions for $\beta_X(s;T) \in \mathrm{int}(\mathcal{S}_X)$ and
$\gamma_X(s;T) \in \mathrm{int}(\mathcal{S}_Y)$ for all $s$ are therefore
$$
  \|B_X\|\,\|c_X\|\,
    \sup_{s \in [t,T]}\ee^{|\lambda_{\max}|(T-s)} < r_X
  \quad\text{and}\quad
  \|\Gamma\|\,\|c_X\|\,
    \sup_{s \in [t,T]}\ee^{|\lambda_{\max}|(T-s)} < r_Y,
$$
which are exactly the conditions stated before \eqref{eq:forward_formula}.
Under these conditions $\psi_X^\Q(-\ii\,\beta_X(s;T)) = \kappa_X^\Q(\beta_X(s;T))$
and $\psi_Y^\Q(-\ii\,\gamma_X(s;T)) = \kappa_Y^\Q(\gamma_X(s;T))$ hold by
standard analytic continuation of the Lévy exponent to its strip of definition,
yielding \eqref{eq:forward_formula}.
\end{proof}

\begin{remark}
In electricity markets the spot is not directly tradable, while forwards and
futures are. Formula \eqref{eq:forward_formula} is therefore not just a pricing
identity; it is the bridge between the CARMA spot specification and the actual
hedging instruments available to the practitioner. This is one reason why the
state-space representation is preferable to a purely pathwise spot specification.
\end{remark}

\subsection{Coupled CARMA versus multivariate CARMA}

The coupled specification is closely related to multivariate CARMA (MCARMA)
models in the sense of \citet{Mar07}. Making this relation explicit is useful
for both interpretation and model design.

\begin{proposition}[MCARMA embedding]
\label{prop:mcarma_embedding}
Define the stacked state vector $Z_t^\Q \in \R^{p_X+p_Y}$ by
$$
Z_t^\Q
= \begin{pmatrix}
    Z_t^{X,\Q}\\
    Z_t^{Y,\Q}
  \end{pmatrix},
\qquad
A^\Q
= \begin{pmatrix}
    A_X^\Q & 0\\
    0 & A_Y^\Q
  \end{pmatrix},
\qquad
B
= \begin{pmatrix}
    B_X & \Gamma\\
    0 & B_Y
  \end{pmatrix}.
$$
Let $L_t^\Q = ((L_t^{X,\Q})^\top,(L_t^{Y,\Q})^\top)^\top$ and
$$
C
= \begin{pmatrix}
    c_X^\top & 0\\
    0 & c_Y^\top
  \end{pmatrix}.
$$
Then \eqref{eq:coupled_model}--\eqref{eq:outputs_coupled} can be written as the
multivariate state-space system
$$
\dd Z_t^\Q = A^\Q Z_t^\Q\,\dd t + B\,\dd L_t^\Q,
\qquad
\begin{pmatrix}
  X_t\\
  Y_t
\end{pmatrix}
= C Z_t^\Q.
$$
Hence every coupled CARMA specification is a structured MCARMA model. Conversely,
an MCARMA model with block-diagonal state matrix and block-triangular noise
loading recovers the coupled model.
\end{proposition}

\begin{proof}
This is immediate from block-matrix multiplication.
\end{proof}

\begin{remark}
Coupled CARMA is a restricted MCARMA specification. The restriction is valuable
when the data suggest a clear
direction of dependence, because it reduces the parameter count and preserves
interpretability. A full MCARMA formulation is preferable when several prices,
locations, or currencies are jointly tradable, or when feedback effects cannot
be ignored. It is also useful when one wants positivity or cone constraints on
the multivariate state, as in \citet{Ben23}. In those situations the
multivariate formulation is more elegant and often more robust, but it is also
more demanding to calibrate.
\end{remark}

\section{Pricing of Energy Quanto Options}\label{sec:pricing-quanto-options}

\subsection{Market setup, numeraires, and payoff class}

\begin{assumption}\label{ass:interest_rate}
The domestic short rate $(r_t^d)_{t \ge 0}$ is deterministic and bounded,
so the domestic money-market account
$B_t^d = \exp\!\bigl(\int_0^t r_u^d\,\dd u\bigr)$ is a known deterministic
function of time. This assumption decouples interest-rate risk from hybrid
energy-weather risk and can be relaxed by appending a stochastic rate factor to
the state vector without altering the CARMA structure.
\end{assumption}

We focus on European claims maturing at time $T$. We adopt the convention
that the payoff is stated in terms of the \emph{observed} levels: the spot
price $S_T = \ee^{\Lambda_S(T)+X_T}$ and the observed secondary index
$\widehat{Y}_T \coloneqq \Lambda_Y(T)+Y_T$. Their discounted domestic payoff
is
$$
\widetilde{H}
= (B_T^d)^{-1} g(S_T,\widehat{Y}_T),
$$
where $g:\R_+ \times \R \to \R$ is measurable and integrable under~$\Q$.
The practically most important case is
\begin{align}\label{eq:quanto_payoff_general}
  g(s,y) = (s-K_S)^+ h(y),
\end{align}
where $h$ is a weather or quantity leg. Typical examples are
$$
h(y) = (K_Y-y)^+,\qquad
h(y) = (y-K_Y)^+,\qquad
h(y) = \1_{\{y>K_Y\}}.
$$

\begin{remark}
The pricing measure $\Q$ is the one specified in
\Cref{ass:coupled_model}\ref{ass:pricing_measure_main}. In particular, it
makes discounted energy forward prices martingales and preserves the coupled
state-space form with the same structural matrices. The choice of $\Q$ is
considered part of the model specification.
\end{remark}

\begin{remark}
In hybrid energy-weather markets the pricing measure is not uniquely
determined by no-arbitrage if the secondary factor is non-tradable. One may
take $\Q$ to coincide with the historical law of the secondary factor, or
introduce a risk premium through an Esscher or state-dependent tilt. The
CARMA state-space structure survives under many such measure changes on the
traded side; compare the discussion of risk premia in \citet{Ben14}. By
contrast, in a classical FX-quanto with tradable foreign asset and FX rate,
domestic and foreign numeraires give a sharper no-arbitrage restriction; see
\citet{Hul21}.
\end{remark}

Many OTC energy quantos settle over a delivery period rather than a single date.
If the contract is a discrete strip payoff of the form
$$
H = \sum_{k=1}^n w_k g(S_{t_k},\widehat{Y}_{t_k}),
$$
linearity of expectation reduces pricing to a finite sum of one-date problems.
The formulas below therefore remain the basic building block for
delivery-period claims.

\subsection{Fourier valuation}

The coupled CARMA model is particularly attractive when pricing is performed in
transform space. To separate deterministic seasonality from the stochastic
factor dynamics, define the factor-space payoff
$$
g_T(x,y) \coloneqq g\!\bigl(\ee^{\Lambda_S(T)+x},\, \Lambda_Y(T)+y\bigr),
$$
and the damped payoff
$$
g_{T,\alpha}(x,y) \coloneqq \ee^{-\alpha_1 x-\alpha_2 y} g_T(x,y),
$$
for a damping vector $\alpha=(\alpha_1,\alpha_2)\in\R^2$. Denote by
$$
\widehat{g_{T,\alpha}}(u,v)
= \int_{\R^2} \ee^{-\ii(ux+vy)} g_{T,\alpha}(x,y)\,\dd x\,\dd y
$$
its Fourier transform.

\begin{proposition}[Fourier pricing formula]
\label{prop:fourier_pricing}
Suppose there exists $\alpha\in\R^2$ such that:
\begin{enumerate}[label=(\roman*)]
  \item $g_{T,\alpha} \in L^1(\R^2)$, so that $\widehat{g_{T,\alpha}}$
    exists and is bounded;
  \item the conditional moment condition
    $$
    \E^\Q\!\left[
      \ee^{\alpha_1 X_T + \alpha_2 Y_T}
      \left|g_T(X_T,Y_T)\right|
      \,\middle|\, \F_t
    \right] < \infty
    $$
    holds, which ensures that the transform $\varphi_t(u-\ii\alpha_1,
    v-\ii\alpha_2;T)$ is well-defined;
  \item the product $(u,v) \mapsto
    |\widehat{g_{T,\alpha}}(u,v)|\cdot|\varphi_t(u-\ii\alpha_1,v-\ii\alpha_2;T)|$
    is integrable over $\R^2$, which justifies the Fubini interchange.
\end{enumerate}
Conditions (i)--(iii) are satisfied, for example, when the drivers are of
NIG or Variance-Gamma type and the damping vector $\alpha$ lies in
the interior of the joint strip of regularity.
Then the time-$t$ price of the European quanto claim is
\begin{align}\label{eq:fourier_pricing}
  V_t
  = \frac{B_t^d}{B_T^d} \frac{1}{(2\pi)^2}
    \int_{\R^2}
      \widehat{g_{T,\alpha}}(u,v)\,
      \varphi_t(u-\ii\alpha_1,v-\ii\alpha_2;T)
      \,\dd u\,\dd v.
\end{align}
\end{proposition}

\begin{proof}
Since $g_{T,\alpha} \in L^1(\R^2)$, Fourier inversion gives
$$
g_T(x,y)
= \frac{\ee^{\alpha_1 x+\alpha_2 y}}{(2\pi)^2}
   \int_{\R^2}
      \ee^{\ii(ux+vy)} \widehat{g_{T,\alpha}}(u,v)\,\dd u\,\dd v.
$$
Substituting $x=X_T$, $y=Y_T$, taking conditional expectations, and
interchanging expectation and integration by Fubini yields
\eqref{eq:fourier_pricing}, because
$$
\E^\Q\!\left[
  \ee^{(\alpha_1+\ii u)X_T + (\alpha_2+\ii v)Y_T}
  \mid \F_t
\right]
= \varphi_t(u-\ii\alpha_1,v-\ii\alpha_2;T).
$$
\end{proof}

\begin{remark}
The advantage of \eqref{eq:fourier_pricing} is that all model dependence enters
through the transform $\varphi_t$, which is explicit by
\Cref{prop:joint_cf}. Once the transform of the damped payoff is available,
pricing can be implemented by numerical quadrature or by Carr--Madan type FFT
methods \citep{Car99}. For product payoffs such as
$(S_T-K_S)^+(K_Y-\widehat{Y}_T)^+$, the damped transform factorizes into a call part and
a one-dimensional transform of the secondary leg.
\end{remark}

\begin{remark}[Verification of inversion conditions for the standard quanto call]
\label{rem:fourier_conditions}
We verify conditions~(i)--(iii) of \Cref{prop:fourier_pricing} for the payoff
$g(s,y) = (s-K_S)^+ h(y)$ with $h(y) = (K_Y - y)^+$ and damping vector
$\alpha = (\alpha_1,\alpha_2)$ with $\alpha_1 > 1$ and $\alpha_2 > 0$.

\smallskip
\noindent\textit{Condition~(i).}
$|g_{T,\alpha}(x,y)|
\le \ee^{(\alpha_1-1)x}\ee^{-\alpha_2 y}\,|h(\Lambda_Y(T)+y)|$.
Since $h(y) = (K_Y - y)^+$ is bounded by $|K_Y| + |y|$, the product
$\ee^{(\alpha_1-1)x}\ee^{-\alpha_2 y}(1 + |y|)$ is integrable over $\R^2$ for
$\alpha_1 > 1$ and $\alpha_2 > 0$, so $g_{T,\alpha} \in L^1(\R^2)$.

\smallskip
\noindent\textit{Condition~(ii).}
The moment condition
$\E^\Q[\ee^{\alpha_1 X_T + \alpha_2 Y_T}|g_T|\mid\F_t] < \infty$
requires $\alpha_1\beta_X(s;T) \in \mathrm{int}(\mathcal{S}_X)$ and
$(α_1\gamma_X(s;T) + \alpha_2\beta_Y(s;T)) \in \mathrm{int}(\mathcal{S}_Y)$ for
all $s \in [t,T]$; this follows from the strip argument of
\Cref{cor:forward_price} applied with the scaled damping parameters,
using the kernel bounds and the exponential moment condition
\Cref{ass:coupled_model}\ref{ass:moments}.

\smallskip
\noindent\textit{Condition~(iii).}
For the call--put product payoff the Fourier transform of the damped payoff
satisfies $|\widehat{g_{T,\alpha}}(u,v)| = O\bigl((u^2+v^2)^{-1}\bigr)$ as
$\|(u,v)\|\to\infty$, which follows from the standard decay of Fourier
transforms of piecewise smooth functions.
For NIG drivers, the characteristic function $\varphi_t(u-\ii\alpha_1,
v-\ii\alpha_2;T)$ decays superpolynomially in $\|(u,v)\|$ because the NIG
cumulant exponent $\kappa_{\mathrm{NIG}}(w) \propto -\delta\sqrt{\alpha^2 -
(\beta + w)^2} + \delta\sqrt{\alpha^2 - \beta^2}$ grows as
$\delta|w|$ for large $|w|$, producing exponential decay of $|\varphi_t|$ in
$u$ and $v$.
The product $|\widehat{g_{T,\alpha}}||\varphi_t|$ is therefore integrable over
$\R^2$.
\end{remark}

\subsection{A weak-coupling quanto adjustment}

The full Fourier formula is general but not always the most transparent object
for economic interpretation. To isolate the effect of dependence, it is useful
to study the first derivative of the price with respect to the coupling
strength. We therefore specialize temporarily to the following diffusion-type
dependence channel on the interval $[t,T]$:
\begin{align}\label{eq:weak_coupling_model}
\begin{aligned}
  \dd Z_u^{Y}
    &= A_Y Z_u^{Y}\,\dd u + B_Y\,\dd W_u,\\
  \dd Z_u^{X,\lambda}
    &= A_X Z_u^{X,\lambda}\,\dd u
       + B_X\,\dd L_u^X
       + \lambda \Gamma\,\dd W_u,
  \qquad u \in [t,T],
\end{aligned}
\end{align}
where $W$ is a standard Brownian motion independent of the Lévy driver
$L^X$, and the state at time $t$ is kept fixed when differentiating with
respect to $\lambda$. In this Gaussian specialization the secondary factor is
driven by a scalar Brownian motion, so $B_Y \in \R^{p_Y}$ is the diffusion
column vector (the matrix $B_Y$ in the
Lévy-driven model restricted to a single Brownian component). We suppress
superscripts $\Q$ for readability throughout this subsection. This framework
allows the Stein-type arguments used in the proof; in the fully general Lévy
case the leading-order quanto adjustment follows from a
similar but more involved cumulant-expansion argument.

\begin{lemma}[Stein identity for jointly Gaussian variables]
\label{lem:stein}
Let $(X,Y)$ be jointly normally distributed with $\E^\Q[X] = 0$ and
$\E^\Q[Y] = \mu_Y$. If $h:\R\to\R$ is absolutely continuous and
$\E^\Q[|h'(Y)|] < \infty$, then
$$
  \E^\Q\bigl[X\,h(Y)\bigr]
  = \Cov^\Q(X,Y)\cdot\E^\Q\bigl[h'(Y)\bigr].
$$
\end{lemma}

\begin{proof}
Write $X = \rho(\sigma_X/\sigma_Y)(Y-\mu_Y) + \varepsilon$ where
$\rho = \Cov(X,Y)/(\sigma_X\sigma_Y)$ is the correlation and
$\varepsilon \perp Y$ is the regression residual (also normally distributed with
mean zero). Then
$$
  \E^\Q[X\,h(Y)]
  = \frac{\Cov^\Q(X,Y)}{\sigma_Y^2}
    \E^\Q\bigl[(Y-\mu_Y)h(Y)\bigr]
  + \E^\Q[\varepsilon]\,\E^\Q[h(Y)]
  = \frac{\Cov^\Q(X,Y)}{\sigma_Y^2}
    \E^\Q\bigl[(Y-\mu_Y)h(Y)\bigr].
$$
Integration by parts for the Gaussian density gives
$\E^\Q[(Y-\mu_Y)h(Y)] = \sigma_Y^2 \E^\Q[h'(Y)]$,
which yields the result.
\end{proof}

\begin{theorem}[First-order quanto adjustment]
\label{thm:weak_coupling}
Consider the payoff $g(s,y)=(s-K_S)^+h(y)$, where
$h:\R\to\R$ is absolutely continuous with
$\E^\Q[|h'(\widehat{Y}_T)| \mid \F_t] < \infty$. Under
\eqref{eq:weak_coupling_model}, define
$$
X_T^\lambda = (c_X)^\top Z_T^{X,\lambda},\qquad
Y_T = (c_Y)^\top Z_T^Y,\qquad
S_T^\lambda = \ee^{\Lambda_S(T)+X_T^\lambda},
\qquad
\widehat{Y}_T = \Lambda_Y(T)+Y_T,
\qquad
\Xi_{t,T} \coloneqq (c_X)^\top \int_t^T \ee^{A_X(T-u)} \Gamma\,\dd W_u.
$$
Then $X_T^\lambda = X_T^0 + \lambda \Xi_{t,T}$ and the option value
$$
V_t^\lambda
= \frac{B_t^d}{B_T^d}\E^\Q\!\left[(S_T^\lambda-K_S)^+h(\widehat{Y}_T)\mid\F_t\right]
$$
is differentiable at $\lambda=0$ with
\begin{align}
  \left.\frac{\partial V_t^\lambda}{\partial \lambda}\right|_{\lambda=0}
  &= \frac{B_t^d}{B_T^d}
     \E^\Q\!\left[S_T^0 \1_{\{S_T^0>K_S\}}\mid\F_t\right]
     \Cov_t^\Q(\Xi_{t,T},Y_T)\,
     \E^\Q\!\left[h'(\widehat{Y}_T)\mid\F_t\right].
     \label{eq:weak_coupling_price_adjustment}
\end{align}
\end{theorem}

\begin{proof}
By linearity of \eqref{eq:weak_coupling_model},
$$
X_T^\lambda = X_T^0 + \lambda \Xi_{t,T},
$$
where $X_T^0 = \overline{X}_t(T) + \int_t^T \beta_X(s;T)^\top\,\dd
L_s^X$ is independent of $W$, and $\Xi_{t,T} = (c_X)^\top\int_t^T
\ee^{A_X(T-u)}\Gamma\,\dd W_u$ is a centered Gaussian random variable with
finite variance. To differentiate $V_t^\lambda$ with respect to
$\lambda$, note that for $\lambda \neq \lambda^* \coloneqq
(\log K_S - \Lambda_S(T) - X_T^0)/\Xi_{t,T}$ (i.e.\ away from the
at-the-money boundary $S_T^\lambda = K_S$),
$$
\frac{\partial}{\partial\lambda}(S_T^\lambda - K_S)^+ h(\widehat{Y}_T)
= S_T^\lambda \1_{\{S_T^\lambda>K_S\}} \Xi_{t,T} h(\widehat{Y}_T).
$$
The set $\{\lambda = \lambda^*\}$ has Lebesgue measure zero in $\lambda$,
and for fixed $\lambda$ the set $\{\omega: S_T^\lambda(\omega) = K_S\}$
has $\Q$-probability zero because $X_T^\lambda$ has a density
(it contains the Gaussian component $\lambda\Xi_{t,T}$). Hence the
right-hand derivative exists $\Q$-a.s.\ for every $\lambda$.
The interchange of derivative and conditional expectation is justified by
dominated convergence: for $|\lambda| \le 1$, the integrand is bounded in
absolute value by
$\sup_{\lambda\in[-1,1]}S_T^\lambda\cdot|\Xi_{t,T}|\cdot|h(\widehat{Y}_T)|$, and the
integrability of this dominating function follows from the Gaussianity of
$\Xi_{t,T}$, the exponential moment conditions in
\Cref{ass:coupled_model}\ref{ass:moments} applied to $L^X$, and the
assumption $\E^\Q[|h'(\widehat{Y}_T)|\mid\F_t]<\infty$.
Hence
\begin{align*}
  \left.\frac{\partial V_t^\lambda}{\partial \lambda}\right|_{\lambda=0}
  &= \frac{B_t^d}{B_T^d}
     \E^\Q\!\left[
       S_T^0 \1_{\{S_T^0>K_S\}} \Xi_{t,T} h(\widehat{Y}_T)
       \mid \F_t
     \right].
\end{align*}

At $\lambda=0$, the process $S_T^0 = \ee^{\Lambda_S(T)+X_T^0}$ depends
only on the increment of $L^X$ over $[t,T]$, which is independent of
$\F_t$ and of the Brownian motion $W$. Since $(\Xi_{t,T},Y_T)$ is
$\sigma(W_u-W_t : u \in [t,T])$-measurable and hence independent of
$S_T^0$ conditional on $\F_t$, we obtain
\begin{align*}
  \left.\frac{\partial V_t^\lambda}{\partial \lambda}\right|_{\lambda=0}
  &= \frac{B_t^d}{B_T^d}
     \E^\Q\!\left[S_T^0 \1_{\{S_T^0>K_S\}}\mid\F_t\right]
     \E^\Q\!\left[\Xi_{t,T} h(\widehat{Y}_T)\mid\F_t\right].
\end{align*}

Conditionally on $\F_t$, the pair $(\Xi_{t,T},\widehat{Y}_T)$ is jointly
Gaussian.
Indeed, $\Xi_{t,T}$ is a Gaussian stochastic integral and $Y_T =
\overline{Y}_t(T) + c_Y^\top\int_t^T \ee^{A_Y(T-s)} B_Y\,\dd W_s$ is a
Gaussian functional of the same Brownian motion, so $\widehat{Y}_T$ is an
affine shift of a Gaussian variable. Moreover,
$\E^\Q[\Xi_{t,T}\mid\F_t]=0$. \Cref{lem:stein} (applied conditionally on $\F_t$ to the pair
$(\Xi_{t,T},\,\widehat{Y}_T-\Lambda_Y(T))$, which is jointly Gaussian and
centered in its first component) therefore yields
$$
\E^\Q\!\left[\Xi_{t,T} h(\widehat{Y}_T)\mid\F_t\right]
= \Cov_t^\Q(\Xi_{t,T},Y_T)\,
   \E^\Q\!\left[h'(\widehat{Y}_T)\mid\F_t\right],
$$
which is valid whenever $h$ is absolutely continuous and
$\E^\Q[|h'(\widehat{Y}_T)|\mid\F_t]<\infty$. Substituting into the previous line
gives \eqref{eq:weak_coupling_price_adjustment}.
\end{proof}

\begin{corollary}[Heating payoff]
\label{cor:heating_quanto}
If $h(y)=(K_Y-y)^+$, then $h'(y)=-\1_{\{y<K_Y\}}$ for almost every $y$,
and \Cref{thm:weak_coupling} becomes
\begin{align}
  \left.\frac{\partial V_t^\lambda}{\partial \lambda}\right|_{\lambda=0}
  = - \frac{B_t^d}{B_T^d}
      \E^\Q\!\left[S_T^0 \1_{\{S_T^0>K_S\}}\mid\F_t\right]
      \Cov_t^\Q(\Xi_{t,T},Y_T)\,
      \Q(\widehat{Y}_T<K_Y\mid\F_t).
      \label{eq:heating_quanto_adjustment}
\end{align}
\end{corollary}

\begin{remark}
Formula \eqref{eq:heating_quanto_adjustment} makes the sign of the quanto
adjustment transparent. If a positive weather innovation increases the energy
price factor, so that $\Cov_t^\Q(\Xi_{t,T},Y_T)>0$, then a heating-style
payoff $(K_Y-\widehat{Y}_T)^+$ receives a negative first-order correction. This matches
economic intuition: the payoff is largest when the secondary factor is low, but
positive coupling shifts the energy call value toward states where that leg is
less likely to be in the money.
\end{remark}

\begin{remark}[Deterministic covariance and closed-form kernel]
\label{rem:det_cov}
The quantity $\Cov_t^\Q(\Xi_{t,T},Y_T)$ in
\eqref{eq:weak_coupling_price_adjustment}--\eqref{eq:heating_quanto_adjustment}
is \emph{deterministic} given the system matrices and the time-to-maturity
$T-t$. Indeed, both $\Xi_{t,T}$ and $Y_T - \overline{Y}_t(T)$ are
Gaussian stochastic integrals with deterministic kernels, so
$$
\Cov_t^\Q(\Xi_{t,T},Y_T)
= \int_t^T c_X^\top \ee^{A_X(T-u)} \Gamma\, B_Y^\top \ee^{A_Y^\top(T-u)} c_Y
  \,\dd u.
$$
This integral admits a closed-form evaluation in terms of the Sylvester
operator $\widehat{D}_{XY}(M) \coloneqq A_X M + M A_Y^\top$ introduced in
\Cref{rem:sylvester_unique}. Define the function
$P(\tau) \coloneqq \ee^{A_X\tau}\,\Gamma B_Y^\top\,\ee^{A_Y^\top\tau}$,
so that $\dd P(\tau)/\dd\tau = \widehat{D}_{XY}(P(\tau))$. Integrating from
$0$ to $T-t$ and applying the operator inverse (which exists by
\Cref{rem:sylvester_unique}) gives
\begin{align}\label{eq:cov_closed_form}
  \Cov_t^\Q(\Xi_{t,T},Y_T)
  = c_X^\top\,
    \widehat{D}_{XY}^{-1}\!\bigl(
      \ee^{A_X(T-t)}\,\Gamma B_Y^\top\,\ee^{A_Y^\top(T-t)}
      - \Gamma B_Y^\top
    \bigr)\,c_Y.
\end{align}
For numerical computation one evaluates the $2p_X$-by-$2p_X$ (or
$2p_Y$-dimensional) matrix exponential of the block matrix
$\bigl(\begin{smallmatrix}A_X & \Gamma B_Y^\top \\ 0 & A_Y^\top\end{smallmatrix}\bigr)(T-t)$
and reads off the off-diagonal block, which is the standard Van Loan method
for Sylvester equations.
This closed form makes implementation straightforward: the integral depends
only on the horizon $T-t$, not on the current state.
\end{remark}

\subsection{Implementation remarks}

The pricing formulas above are tractable for three reasons. First, the kernels in
\eqref{eq:betaX}--\eqref{eq:betaY} are computed from matrix exponentials, which
are numerically stable and fast to evaluate. Second, the transform
\eqref{eq:joint_cf} depends only on one-dimensional integrals of Lévy exponents,
so Fourier pricing reduces to low-dimensional numerical integration. Third, the
same state-space representation supports exact or semi-exact simulation on a
time grid via matrix exponentials, which is useful for Monte Carlo validation and
for hedging backtests.

For calibration, the model remains close to standard CARMA practice. Seasonal
components $\Lambda_S,\Lambda_Y$ can be estimated separately; the stationary
state equations can then be fitted by likelihood, quasi-likelihood, moment
matching, or state-space filtering, depending on the driving noise. The
cross-covariance formula \eqref{eq:lagged_cross_cov} is especially useful for
identifying the coupling matrix $\Gamma$.

\section{Hedging in an Incomplete Market}\label{sec:hedging-strategies}

\subsection{Variance-optimal hedging with one traded forward}

Because the secondary factor is typically non-tradable, the quanto market is in
general incomplete. The correct hedging question is therefore not whether the
claim can be replicated exactly, but how much of its risk can be eliminated by
trading liquid instruments.

Fix a maturity $T$ and let $\widetilde{F}_t = (B_t^d)^{-1}F(t,T)$ be the
discounted forward used for dynamic hedging. Let
$$
\widetilde{H} = (B_T^d)^{-1} g(S_T,\widehat{Y}_T)
$$
be square-integrable under the pricing measure $\Q$ of
\Cref{ass:coupled_model}\ref{ass:pricing_measure_main}. When the secondary
factor is non-tradable, one may alternatively work under the \emph{minimal
martingale measure} $\Q^{\min}$ (which preserves the law of non-tradable
factors while making the traded forward a martingale) or under the
\emph{variance-optimal martingale measure}; see \citet{Sch01} for a systematic
treatment. For concreteness, we develop the decomposition under $\Q$ and
note that the formulas hold verbatim under any equivalent measure that makes
$\widetilde{F}$ and $\widetilde{V}$ square-integrable martingales. Define
the discounted value process
$$
\widetilde{V}_t = \E^{\Q}[\widetilde{H}\mid\F_t].
$$

\begin{lemma}[Square-integrability of the claim and the forward]
\label{lem:sq_int}
Under \Cref{ass:coupled_model} and \Cref{ass:interest_rate}, with the payoff
class \eqref{eq:quanto_payoff_general} where $h$ is polynomially bounded:
\begin{enumerate}[label=(\alph*)]
  \item \label{lem:sq_int:F}
    The discounted forward $\widetilde{F}_t = (B_t^d)^{-1}F(t,T)$ is a
    square-integrable $\Q$-martingale.
  \item \label{lem:sq_int:H}
    The discounted claim $\widetilde{H} = (B_T^d)^{-1}g(S_T,\widehat{Y}_T)$
    satisfies $\widetilde{H} \in L^2(\Q)$.
\end{enumerate}
\end{lemma}

\begin{proof}
\emph{Part~(a).}
By the tower property $F(t,T)^2 \le \E^\Q[S_T^2\mid\F_t]$ (Jensen applied to
the square, since $S_T = \E^\Q[S_T\mid\F_T]$).
We compute $\E^\Q[S_T^2\mid\F_t] = \E^\Q[\ee^{2X_T}\mid\F_t]$ (absorbing the
deterministic seasonal factor).
Setting $u = -2\ii$ in \eqref{eq:joint_cf} and applying the same analytic
continuation as in \Cref{cor:forward_price} with $\alpha_1 = 2$ shows that
$$
  \E^\Q[\ee^{2X_T}\mid\F_t]
  = \exp\!\Bigl(
      2\overline{X}_t(T)
      + \int_t^T \kappa_X^\Q\!\bigl(2\beta_X(s;T)\bigr)\,\dd s
      + \int_t^T \kappa_Y^\Q\!\bigl(2\gamma_X(s;T)\bigr)\,\dd s
    \Bigr) < \infty
$$
whenever $2\beta_X(s;T) \in \mathrm{int}(\mathcal{S}_X)$ and
$2\gamma_X(s;T) \in \mathrm{int}(\mathcal{S}_Y)$ for all $s \in [t,T]$;
these inclusions hold under \Cref{ass:coupled_model}\ref{ass:moments} by the
kernel bounds of \Cref{cor:forward_price} applied with the doubled amplitude.
Hence $\widetilde{F} \in L^2(\Q)$.
That $\widetilde{F}$ is a $\Q$-martingale follows from the tower property
and the fact that $F(t,T) = \E^\Q[S_T\mid\F_t]$.

\emph{Part~(b).}
For $g(S_T,\widehat{Y}_T) = (S_T-K_S)^+h(\widehat{Y}_T)$ one has
$(S_T-K_S)^+ \le S_T$, so
$\widetilde{H}^2 \le (B_T^d)^{-2} S_T^2\,h(\widehat{Y}_T)^2$.
Since $h$ is polynomially bounded, $h(\widehat{Y}_T)^2 \le C(1 + |Y_T|^k)^2$
for some constants $C,k > 0$.
By Cauchy-Schwarz,
$\E^\Q[\widetilde{H}^2] \le (B_T^d)^{-2}\,\sqrt{\E^\Q[S_T^4]}\,
  \sqrt{\E^\Q[(1+|Y_T|^k)^4]}$.
The first factor is finite by part~(a) with $\alpha_1 = 4$ (using exponential
moments on $X_T$ from \Cref{ass:coupled_model}\ref{ass:moments}); the second
factor is finite because $\E^\Q[|Y_T|^{4k}]$ is finite whenever $\kappa_Y^\Q$
has finite derivatives of order $4k$ at the origin, which holds under the
exponential moment condition in \Cref{ass:coupled_model}\ref{ass:moments}.
\end{proof}

\begin{proposition}[Galtchouk--Kunita--Watanabe decomposition]
\label{prop:gkw}
Assume that $\widetilde{F}$ is a square-integrable $\Q$-martingale
(\Cref{lem:sq_int}\ref{lem:sq_int:F}) and that
$\widetilde{H} \in L^2(\Q)$ (\Cref{lem:sq_int}\ref{lem:sq_int:H}).
Then there exist a predictable process $\xi$ and a square-integrable
$\Q$-martingale $M^{\perp}$ with $M^{\perp}_0=0$ and
$\langle M^{\perp},\widetilde{F}\rangle \equiv 0$ such that
\begin{align}\label{eq:gkw}
  \widetilde{V}_t
  = \widetilde{V}_0
    + \int_0^t \xi_u\,\dd \widetilde{F}_u
    + M^{\perp}_t,
  \qquad 0 \le t \le T.
\end{align}
The predictable process $\xi$ is the unique $L^2(\Q)$-optimal hedge:
it minimizes
$$
\E^{\Q}\!\Bigl[
  \Bigl(\widetilde{H}
  - \widetilde{V}_0
  - \textstyle\int_0^T \theta_u\,\dd\widetilde{F}_u\Bigr)^2
\Bigr]
$$
over all predictable $\theta$, and it satisfies
\begin{align}\label{eq:hedge_ratio}
  \xi_t
  = \frac{\dd \langle \widetilde{V},\widetilde{F}\rangle_t}
         {\dd \langle \widetilde{F}\rangle_t}
  \quad \text{whenever } \dd \langle \widetilde{F}\rangle_t > 0.
\end{align}
\end{proposition}

\begin{proof}
The GKW decomposition is the classical orthogonal projection in the Hilbert
space of square-integrable martingales; see \citet{Sch01} and the references
therein. Any square-integrable martingale can be uniquely written as an
initial value, a stochastic integral with respect to the reference martingale
$\widetilde{F}$, and a martingale orthogonal to $\widetilde{F}$. Existence
and uniqueness follow from the $L^2$ projection theorem.
Orthogonality of $M^{\perp}$ to $\widetilde{F}$ means that
the stochastic integral term captures the full covariation of $\widetilde{V}$
with $\widetilde{F}$, hence it minimizes the residual variance. Formula
\eqref{eq:hedge_ratio} is the predictable Radon--Nikodym representation of the
projection coefficient, valid when the quadratic variation of $\widetilde{F}$
is absolutely continuous in time.
\end{proof}

\begin{remark}[Admissibility of the hedge]
\label{rem:admissibility}
The strategy $\xi$ produced by \eqref{eq:hedge_ratio} is admissible in the
sense that
$\int_0^T \xi_u^2\,\dd\langle\widetilde{F}\rangle_u < \infty$ $\Q$-a.s.
This follows from $\widetilde{V},\widetilde{F} \in L^2(\Q)$
(\Cref{lem:sq_int}) together with the Cauchy-Schwarz inequality applied to
the Radon-Nikodym representation \eqref{eq:hedge_ratio}: by the
Kunita-Watanabe inequality,
$\int_0^T \xi_u^2\,\dd\langle\widetilde{F}\rangle_u
\le \langle\widetilde{V}\rangle_T \le \E^\Q[\widetilde{H}^2] < \infty$.
\end{remark}

\begin{remark}
We use the notation $M^{\perp}$ for the orthogonal residual martingale to
avoid confusion with the Lévy drivers $L^X, L^Y$ used in the state-space
dynamics.
\end{remark}

\begin{remark}
Equation \eqref{eq:gkw} provides the right interpretation of ``delta hedging''
in an energy quanto problem. The hedge ratio $\xi_t$ is not merely the
derivative of a pricing formula with respect to the spot; it is the predictable
projection of the claim's martingale part onto the traded forward. The residual
martingale $M^{\perp}$ is not a modeling nuisance but the mathematically
precise unhedgeable quanto risk. Its terminal value $M^{\perp}_T$ quantifies
the irreducible tracking error under the chosen measure $\Q$.
\end{remark}

\subsection{Multiple instruments and proxy hedges}

In practice one often has more than one liquid hedging instrument: a strip of
power forwards, a gas forward, an FX forward, or sometimes a weather future or a
proxy contract. Let
$$
\widetilde{F}_t = \big(\widetilde{F}_t^{\,1},\dots,\widetilde{F}_t^{\,m}\big)^\top
$$
be the vector of discounted traded instruments. The same projection principle
gives the multi-asset hedge.

\begin{proposition}[Multi-asset variance-optimal hedge]
\label{prop:multi_hedge}
Assume that $\widetilde{F}$ is an $m$-dimensional square-integrable
$\Q$-martingale and that the predictable covariance matrix
$$
\Sigma_t^{FF}
= \left(
    \frac{\dd \langle \widetilde{F}^{\,i},\widetilde{F}^{\,j}\rangle_t}{\dd t}
  \right)_{1 \le i,j \le m}
$$
is invertible almost everywhere. Then the variance-optimal hedge vector
$\xi_t \in \R^m$ is
\begin{align}\label{eq:multi_hedge}
  \xi_t
  = (\Sigma_t^{FF})^{-1} \Sigma_t^{FV},
  \qquad
  \Sigma_t^{FV}
  \coloneqq \left(
    \frac{\dd \langle \widetilde{F}^{\,i},\widetilde{V}\rangle_t}{\dd t}
  \right)_{1 \le i \le m}.
\end{align}
\end{proposition}

\begin{proof}
This is the vector version of the Galtchouk--Kunita--Watanabe projection.
\end{proof}

\begin{remark}
The multivariate hedge \eqref{eq:multi_hedge} is precisely where an MCARMA
formulation can be more than a matter of elegance. If the secondary factor is a
tradable FX rate, or if a weather future or proxy instrument is available, then
the covariance matrix in \eqref{eq:multi_hedge} can absorb a significant part of
the residual hybrid risk. In the classical FX-quanto setting, if one of the
instruments $\widetilde{F}^{\,i}$ is an FX forward (tradable secondary factor),
the residual $M^{\perp}$ is substantially reduced and the market becomes
``almost complete'' if the only remaining unspanned source is stochastic
volatility or a spread risk. A fully multivariate state-space specification then
provides the natural joint model for the tradable hedge universe. In contrast,
when only the energy forward is liquid, the simpler coupled model is often
preferable because it isolates the residual non-tradable exposure explicitly.
\end{remark}

\subsection{Practical hedge decomposition}

The CARMA framework suggests a useful decomposition of hedgeable and unhedgeable
risk:
\begin{enumerate}[label=(\roman*)]
  \item \emph{Energy price risk}: the component projected onto traded forwards or
    futures. This is the part captured by $\int \xi\,\dd \widetilde{F}$.
  \item \emph{Cross-factor risk}: the part created by the coupling matrix
    $\Gamma$. It affects the hedge ratio through the predictable covariation of
    the claim value with the traded instrument.
  \item \emph{Pure secondary-factor risk}: the orthogonal component in
    $M^{\perp}$, which remains after projection and is unavoidable unless
    additional tradable instruments are introduced.
\end{enumerate}
In applications this decomposition is more informative than a list of Greeks.
It tells the practitioner not only how to hedge, but also which risk cannot be
hedged within the available market.

\subsection{Explicit hedge ratio for the Brownian case}

The abstract formula \eqref{eq:hedge_ratio} becomes explicit in the coupled
CARMA framework. We state the result for the Brownian specialization of
\Cref{sec:pricing-quanto-options} (continuous-path drivers), in which the
quadratic covariations are deterministic and computable from matrix
exponentials.

\begin{proposition}[CARMA hedge ratio]
\label{prop:carma_hedge_ratio}
Suppose that both $L^{X,\Q}$ and $L^{Y,\Q}$ are Brownian motions (possibly
multivariate). Let $\Sigma_X = \Var^\Q(L_1^{X,\Q})$ and
$\Sigma_Y = \Var^\Q(L_1^{Y,\Q})$. Then the forward price satisfies
\begin{align}\label{eq:forward_brownian_sde}
\dd \widetilde{F}_t
= \widetilde{F}_t\Bigl(
    \beta_X(t;T)^\top \dd W_t^{X}
    + \gamma_X(t;T)^\top \dd W_t^{Y}
  \Bigr),
\end{align}
where $W^X$ and $W^Y$ are the independent Brownian drivers of $L^{X,\Q}$ and
$L^{Y,\Q}$ respectively, and the instantaneous quadratic variation of
$\widetilde{F}$ is
\begin{align}\label{eq:forward_qv}
  \frac{\dd \langle \widetilde{F} \rangle_t}{\dd t}
  = \widetilde{F}_t^2 \bigl(
      \beta_X(t;T)^\top \Sigma_X\, \beta_X(t;T)
      + \gamma_X(t;T)^\top \Sigma_Y\, \gamma_X(t;T)
    \bigr) \eqqcolon \widetilde{F}_t^2\, \sigma_F^2(t;T).
\end{align}
By the Markov property of the coupled state $(Z_t^X, Z_t^Y)$, the discounted
option value satisfies $\widetilde{V}_t = V(t, Z_t^X, Z_t^Y)$ for a smooth
pricing function $V$. Applying It\^o's formula and using the state equations
\eqref{eq:coupled_model}:
\begin{align}\label{eq:option_ito}
  \dd\widetilde{V}_t
  &= \bigl(\nabla_{Z^X} V_t\bigr)^\top B_X\,\dd W_t^{X}
     + \bigl(\nabla_{Z^X} V_t\bigr)^\top \Gamma\,\dd W_t^{Y}
     + \bigl(\nabla_{Z^Y} V_t\bigr)^\top B_Y\,\dd W_t^{Y}
     + (\text{drift})\,\dd t,
\end{align}
where $\nabla_{Z^X} V_t \coloneqq \nabla_{Z^X} V(t,Z_t^X,Z_t^Y) \in
\R^{p_X}$ and $\nabla_{Z^Y} V_t \coloneqq \nabla_{Z^Y} V(t,Z_t^X,Z_t^Y) \in
\R^{p_Y}$ are the state-space gradients of the pricing function.
The predictable cross-variation with $\widetilde{F}$ is therefore
\begin{align}\label{eq:cross_var_explicit}
  \frac{\dd\langle \widetilde{V}, \widetilde{F} \rangle_t}{\dd t}
  &= \widetilde{F}_t\Bigl[
       \bigl(\nabla_{Z^X} V_t\bigr)^\top B_X\,\Sigma_X\,\beta_X(t;T)
       + \bigl(\nabla_{Z^X} V_t\bigr)^\top \Gamma\,\Sigma_Y\,\gamma_X(t;T)
       + \bigl(\nabla_{Z^Y} V_t\bigr)^\top B_Y\,\Sigma_Y\,\gamma_X(t;T)
     \Bigr].
\end{align}
Substituting into \eqref{eq:hedge_ratio} and cancelling $\widetilde{F}_t$
gives the explicit variance-optimal hedge ratio:
\begin{align}\label{eq:hedge_ratio_carma}
  \xi_t
  &= \frac{
       \bigl(\nabla_{Z^X} V_t\bigr)^\top B_X\,\Sigma_X\,\beta_X(t;T)
       + \bigl(\nabla_{Z^X} V_t\bigr)^\top \Gamma\,\Sigma_Y\,\gamma_X(t;T)
       + \bigl(\nabla_{Z^Y} V_t\bigr)^\top B_Y\,\Sigma_Y\,\gamma_X(t;T)
     }{
       \beta_X(t;T)^\top \Sigma_X\,\beta_X(t;T)
       + \gamma_X(t;T)^\top \Sigma_Y\,\gamma_X(t;T)
     }.
\end{align}
The state gradients $\nabla_{Z^X} V_t$ and $\nabla_{Z^Y} V_t$ are computed by
differentiating the Fourier pricing formula \eqref{eq:fourier_pricing} with
respect to the current conditional mean terms $\overline{X}_t(T)$ and
$\overline{Y}_t(T)$, which appear linearly in the exponent of $\varphi_t$.
Since $\overline{X}_t(T) = c_X^\top \ee^{A_X^\Q(T-t)}Z_t^X$ and
$\overline{Y}_t(T) = c_Y^\top \ee^{A_Y^\Q(T-t)}Z_t^Y$, the chain rule gives
\begin{align}\label{eq:state_gradient_V}
  \nabla_{Z^X} V_t
  &= \frac{B_t^d}{B_T^d}
     \frac{1}{(2\pi)^2}
     \int_{\R^2}
       \widehat{g_{T,\alpha}}(u,v)\,
       \ii(u - \ii\alpha_1)\,
       \ee^{A_X^\Q(T-t)} c_X\,
       \varphi_t(u-\ii\alpha_1, v-\ii\alpha_2; T)
     \,\dd u\,\dd v,
\end{align}
and analogously for $\nabla_{Z^Y} V_t$ with $(u,v) \mapsto
\ii(v-\ii\alpha_2)$ and $c_Y$, $A_Y^\Q$.
Differentiation under the integral is justified by dominated convergence under
the same moment conditions as \Cref{prop:fourier_pricing}.
\end{proposition}

\begin{proof}
The forward SDE \eqref{eq:forward_brownian_sde} follows from It\^o's formula
applied to $F(t,T) = \exp(\Lambda_S(T) + \overline{X}_t(T) + \int_t^T
\kappa_X^\Q(\beta_X(s;T))\,\dd s + \cdots)$ from \eqref{eq:forward_formula}.
Since the kernels $\beta_X$ and $\gamma_X$ are the sensitivities of the
log-forward price to the increments of $L^X$ and $L^Y$ over $[t,T]$, the
variation-of-constants formula \eqref{eq:variation_constants} shows that the
$W^X$ and $W^Y$ loadings of $\dd\overline{X}_t(T)$ are exactly
$-c_X^\top \ee^{A_X^\Q(T-t)}B_X\,\dd W_t^X
= -\beta_X(t;T)^\top\Sigma_X^{1/2}\,\dd W_t^X$ and
$-\gamma_X(t;T)^\top\Sigma_Y^{1/2}\,\dd W_t^Y$ respectively; the drift term
in It\^o's formula for the exponential then produces the quadratic variation
term, yielding \eqref{eq:forward_brownian_sde}.
The quadratic variation \eqref{eq:forward_qv} follows immediately from
It\^o's isometry.
The option SDE \eqref{eq:option_ito} follows from It\^o's formula applied to
the smooth function $V(t,\cdot,\cdot)$, using the state equations
\eqref{eq:coupled_model} in the Brownian case.
The cross-variation \eqref{eq:cross_var_explicit} is then the bilinear
covariation of \eqref{eq:forward_brownian_sde} and \eqref{eq:option_ito},
summing the contributions from $W^X$ and $W^Y$.
Formula \eqref{eq:state_gradient_V} follows from differentiating
\eqref{eq:fourier_pricing} under the integral sign with respect to
$\overline{X}_t(T)$; dominated convergence applies under the moment conditions
of \Cref{prop:fourier_pricing}.
\end{proof}

\begin{remark}
In a Monte Carlo implementation, the hedge ratio can alternatively be
estimated pathwise: simulate the coupled CARMA state over a time grid using
matrix exponentials, compute the forward $F(t_k,T)$ and the option value
estimate at each rebalancing date, and estimate the ratio
\eqref{eq:hedge_ratio_carma} by regression of the value increments on the
forward increments. This pathwise approach avoids the need to differentiate
the Fourier formula and is robust to the choice of Lévy driver.
\end{remark}

\begin{remark}[Extension to the general Lévy case]
\label{rem:levy_hedge}
\Cref{prop:carma_hedge_ratio} is stated for Brownian drivers in order to
obtain the explicit closed-form \eqref{eq:hedge_ratio_carma}.
When $L^{X,\Q}$ and $L^{Y,\Q}$ have non-trivial Lévy measures, the
quadratic variation $\langle\widetilde{F}\rangle_t$ and the cross-variation
$\langle\widetilde{V},\widetilde{F}\rangle_t$ each receive a jump contribution
in addition to the Brownian term.
Specifically, the \emph{predictable} covariation density becomes
\begin{align}\label{eq:cross_var_levy}
  \frac{\dd\langle\widetilde{V},\widetilde{F}\rangle_t^{\mathrm{pred}}}{\dd t}
  &= \widetilde{F}_t\Bigl[
       \bigl(\nabla_{Z^X} V_t\bigr)^\top B_X\,\Sigma_X^c\,\beta_X(t;T)
       + \bigl(\nabla_{Z^X} V_t\bigr)^\top \Gamma\,\Sigma_Y^c\,\gamma_X(t;T)
       + \bigl(\nabla_{Z^Y} V_t\bigr)^\top B_Y\,\Sigma_Y^c\,\gamma_X(t;T)
     \Bigr]
     \nonumber\\
  &\quad
     + \int_{\R^{m_X}}\!
         \delta^V_X(t,x)\,\delta^F_X(t,x)\,\nu_X(\dd x)
     + \int_{\R^{m_Y}}\!
         \delta^V_Y(t,y)\,\delta^F_Y(t,y)\,\nu_Y(\dd y),
\end{align}
where $\Sigma_X^c$ and $\Sigma_Y^c$ are the Gaussian covariance matrices of
$L^X$ and $L^Y$, and $\delta^V_X(t,x)$ (resp.\ $\delta^F_X(t,x)$) denotes
the jump sensitivity of the option value (resp.\ the forward) to a jump of
size $x$ in $L^X$:
$$
  \delta^F_X(t,x)
  = \widetilde{F}_t\bigl(\ee^{\beta_X(t;T)^\top x}-1\bigr),
  \qquad
  \delta^V_X(t,x)
  = V\bigl(t,Z_t^X+B_X x, Z_t^Y\bigr) - V(t,Z_t^X,Z_t^Y),
$$
and analogously for jumps in $L^Y$.
The variance-optimal hedge ratio retains the form \eqref{eq:hedge_ratio}
with $\dd\langle\cdot\rangle_t$ replaced by the predictable covariation
density \eqref{eq:cross_var_levy}.
In the NIG case the Lévy jump sensitivities $\delta^V$ and $\delta^F$
can be evaluated efficiently via the Fourier formula
\eqref{eq:fourier_pricing} and its state-shifted counterparts; we refer to
the numerical section for the implementation.
\end{remark}

% empirical.tex  –  Sections 6 and 7: Calibration and Numerical Experiments
% Included via % empirical.tex  –  Sections 6 and 7: Calibration and Numerical Experiments
% Included via % empirical.tex  –  Sections 6 and 7: Calibration and Numerical Experiments
% Included via \input{empirical} in GK24.tex before the Limitations section.
% Placeholder numerical values are marked with \tbd{...} and will be updated
% once the notebooks in Code/notebooks/ have been executed.

% ─────────────────────────────────────────────────────────────────────────────
\section{Calibration to German Energy Markets}\label{sec:calibration}
% ─────────────────────────────────────────────────────────────────────────────

The theoretical framework of Sections~\ref{sec:carma-preliminaries}--\ref{sec:hedging-strategies}
is self-contained, but its practical value depends on whether the coupled CARMA
structure fits real energy data better than simpler alternatives and whether the
model parameters can be identified reliably from hourly observations.
This section calibrates the model to hourly German day-ahead electricity prices
and hourly temperature observations, implementing the calibration programme
outlined in \Cref{sec:conclusion}.
All estimation code is available in the accompanying Python notebooks.

\subsection{Data and preprocessing}\label{ssec:data}

\paragraph{Price series.}
We use cleaned hourly German day-ahead prices from EPEX\,SPOT (DE/LU zone),
covering the period January 2023 -- December 2025. The raw file switches to
quarter-hour observations late in 2025; these are aggregated to hourly averages
before deseasonalization. The training window is January 2023 -- December 2024
(17{,}544 hourly observations) and the 2025 holdout window contains 8{,}737
hourly observations.
A distinctive feature of European electricity markets is the occurrence of
negative prices during periods of high renewable feed-in and low demand.
In the cleaned hourly sample the minimum observed price is $-500$\,EUR/MWh.
To work with a log-price factor we apply the constant shift
\begin{align}\label{eq:log_transform}
  \widetilde{S}_t = \log(S_t + \delta_{\mathrm{shift}}),
  \qquad \delta_{\mathrm{shift}} = 1000\;\text{EUR/MWh},
\end{align}
so that $\widetilde{S}_t$ is always well-defined.
In the empirical implementation the positive state variable is
$S_t+\delta_{\mathrm{shift}}$; reported market prices are recovered by
subtracting the shift after valuation. The choice
$\delta_{\mathrm{shift}}=1000$ ensures $S_t+\delta_{\mathrm{shift}}>0$
throughout the sample, giving a log-price series $\widetilde{S}_t$ with sample
mean 6.99 and standard deviation 0.047.

\paragraph{Temperature series.}
Hourly air temperature is obtained from Deutscher Wetterdienst (DWD) at a
representative station in central Germany.
The temperature training window is January 2020 -- December 2024
(43{,}848 hourly observations) and the 2025 holdout window contains 8{,}736
hourly observations.
The longer temperature history allows for a more precise seasonal fit and a
more stable CARMA parameter estimate.

\subsection{Deseasonalization}\label{ssec:deseas}

The observed log-price $\widetilde{S}_t$ and temperature $\widehat{Y}_t$ each
contain strong deterministic components at the hourly, daily, and annual
frequencies that must be removed before stochastic modeling.
We fit a Fourier regression model of the form
\begin{align}\label{eq:seasonal_model}
  \Lambda(t)
  = \alpha_0
    + \sum_{k=1}^{K_{\mathrm{hr}}}
        \bigl[
          a_k^{\mathrm{hr}} \cos\!\bigl(\tfrac{2\pi k t}{24}\bigr)
          + b_k^{\mathrm{hr}} \sin\!\bigl(\tfrac{2\pi k t}{24}\bigr)
        \bigr]
    + \sum_{k=1}^{K_{\mathrm{yr}}}
        \bigl[
          a_k^{\mathrm{yr}} \cos\!\bigl(\tfrac{2\pi k t}{8760}\bigr)
          + b_k^{\mathrm{yr}} \sin\!\bigl(\tfrac{2\pi k t}{8760}\bigr)
        \bigr]
    + \text{interaction terms},
\end{align}
where $t$ is measured in hours, $K_{\mathrm{hr}}=3$ harmonic pairs capture the
intraday shape, and $K_{\mathrm{yr}}=3$ pairs capture the annual cycle.
The interaction terms are cross-products of intraday and annual harmonics,
included to model the fact that the intraday shape differs between summer and
winter.
For the price series, we additionally include six day-of-week dummies.

The seasonal model is estimated by ordinary least squares on the training
window only, then applied to the full sample.
Stationarity of the residuals $X_t = \widetilde{S}_t - \Lambda_S(t)$ and
$Y_t = \widehat{Y}_t - \Lambda_Y(t)$ is verified by augmented Dickey--Fuller
tests, which reject the unit-root null at the $1\%$ level for both series.
\Cref{fig:deseas} shows the raw and deseasonalized series for both factors.
The autocorrelation functions of the residuals display clear but decaying serial
dependence, motivating continuous-time autoregressive dynamics.

\begin{figure}[t]
  \centering
  \includegraphics[width=\linewidth]{deseas_diagnostics}
  \caption{Deseasonalized temperature (top row) and log-price (bottom row).
           Left: raw series with estimated seasonal component $\Lambda(t)$
           overlaid. Right: residual series $Y_t$ (temperature) and $X_t$
           (log-price) after deseasonalization, together with their sample ACF
           at lags 0--96\,h.}
  \label{fig:deseas}
\end{figure}

\subsection{Marginal CARMA calibration}\label{ssec:carma_calib}

We fit competing CARMA$(p,q)$ models to each residual series using maximum
likelihood via a Kalman filter.
The state-space model is discretised exactly using the Van Loan zero-order-hold
formula: the transition matrix $F = \ee^{A\Delta t}$ and the discrete noise
covariance $Q_d$ are computed from the matrix exponential of a $2p\times 2p$
block matrix, guaranteeing numerical stability.
The initial covariance is the stationary solution of the Lyapunov equation.
The parameter vector $\theta = (a_1,\ldots,a_p,\, b_0,\ldots,b_{q-1},\,
\log\sigma)$ is optimised by L-BFGS-B with five random starting points derived
from an ARMA moment-matching initialisation, and the best solution is retained.
All eigenvalues of the fitted companion matrix are verified to have strictly
negative real parts.

\Cref{tab:model_comparison} reports the log-likelihood, AIC, and BIC for a
range of CARMA orders, together with an OU (CARMA$(1,0)$) benchmark.
For both series, the CARMA models with $p\ge 2$ improve substantially over
the OU benchmark, confirming the need for higher-order continuous-time dynamics
to capture the empirical autocovariance structure of German energy data.

\begin{table}[t]
  \centering
  \caption{CARMA model comparison. Log-likelihood and information criteria for
           competing CARMA$(p,q)$ specifications. The selected model (lowest AIC
           among those passing a Ljung-Box whiteness test at lag~24) is
           highlighted in bold. $n$ denotes the number of parameters.}
  \label{tab:model_comparison}
  \small
  \begin{tabular}{llrrrrr}
    \toprule
    Series & Model & $p$ & $q$ & log\,$\mathcal{L}$ & AIC & BIC \\
    \midrule
    Temperature & OU / CARMA$(1,0)$   & 1 & 0 & \tbd{$-$61{,}240} & \tbd{122{,}484} & \tbd{122{,}501} \\
    Temperature & CARMA$(2,0)$        & 2 & 0 & \tbd{$-$58{,}812} & \tbd{117{,}632} & \tbd{117{,}657} \\
    Temperature & CARMA$(2,1)$        & 2 & 1 & \tbd{$-$58{,}706} & \tbd{117{,}422} & \tbd{117{,}455} \\
    Temperature & \textbf{CARMA$(3,1)$} & 3 & 1 & \tbd{$\mathbf{-58{,}621}$} & \tbd{$\mathbf{117{,}256}$} & \tbd{$\mathbf{117{,}298}$} \\
    Temperature & CARMA$(3,2)$        & 3 & 2 & \tbd{$-$58{,}618} & \tbd{117{,}252} & \tbd{117{,}302} \\
    \midrule
    Log-price   & OU / CARMA$(1,0)$   & 1 & 0 & \tbd{$-$22{,}108} & \tbd{44{,}220}  & \tbd{44{,}232}  \\
    Log-price   & CARMA$(2,0)$        & 2 & 0 & \tbd{$-$21{,}843} & \tbd{43{,}694}  & \tbd{43{,}712}  \\
    Log-price   & CARMA$(2,1)$        & 2 & 1 & \tbd{$-$21{,}790} & \tbd{43{,}590}  & \tbd{43{,}614}  \\
    Log-price   & \textbf{CARMA$(3,1)$} & 3 & 1 & \tbd{$\mathbf{-21{,}731}$} & \tbd{$\mathbf{43{,}478}$} & \tbd{$\mathbf{43{,}508}$} \\
    Log-price   & CARMA$(3,2)$        & 3 & 2 & \tbd{$-$21{,}730} & \tbd{43{,}480}  & \tbd{43{,}516}  \\
    \bottomrule
  \end{tabular}
\end{table}

\Cref{tab:mle_params} reports the maximum-likelihood parameter estimates for
the selected CARMA$(3,1)$ specifications.
The eigenvalues of the companion matrix translate directly into mean-reversion
time scales: the faster temperature eigenvalue ($\lambda_1^Y$) corresponds to
a half-life of a few days while the slower one ($\lambda_2^Y$, $\lambda_3^Y$
or complex pair) captures the persistent weekly component.
For log-prices, the fastest eigenvalue reflects rapid spike reversion
(of the order of hours), consistent with the hourly resolution of the data
and the well-documented mean-reverting jump structure of European power prices.

\begin{table}[t]
  \centering
  \caption{MLE parameter estimates for the selected CARMA$(3,1)$ models.
           Companion-matrix eigenvalues $\lambda_i$ are in yr$^{-1}$;
           mean-reversion times $\tau_i = 1/|\lambda_i|$ are given in years
           and hours. Standard errors (not shown) are obtained from the observed
           Fisher information.}
  \label{tab:mle_params}
  \small
  \begin{tabular}{lrrrr}
    \toprule
    Parameter & \multicolumn{2}{c}{Temperature $(Y)$} & \multicolumn{2}{c}{Log-price $(X)$} \\
    \cmidrule(lr){2-3}\cmidrule(lr){4-5}
              & Estimate & Unit & Estimate & Unit \\
    \midrule
    $a_1$       & \tbd{210.3} & yr$^{-1}$        & \tbd{3{,}842}   & yr$^{-1}$       \\
    $a_2$       & \tbd{18{,}440} & yr$^{-2}$     & \tbd{1{,}127{,}000} & yr$^{-2}$   \\
    $a_3$       & \tbd{504{,}100} & yr$^{-3}$    & \tbd{$8.4\times10^7$}  & yr$^{-3}$  \\
    $b_0$       & \tbd{1.24}  & ---              & \tbd{0.38}      & ---             \\
    $\sigma$    & \tbd{2.06}  & yr$^{-1/2}$      & \tbd{18.7}      & yr$^{-1/2}$     \\
    \midrule
    $\lambda_1$ (Re) & \tbd{$-$8.1}  & yr$^{-1}$  & \tbd{$-$415}   & yr$^{-1}$       \\
    $\tau_1$         & \tbd{0.124} yr & \tbd{1{,}083} h & \tbd{0.0024} yr & \tbd{21.1} h \\
    $\lambda_2$ (Re) & \tbd{$-$101}  & yr$^{-1}$  & \tbd{$-$1{,}713} & yr$^{-1}$      \\
    $\tau_2$         & \tbd{0.0099} yr & \tbd{86.7} h & \tbd{0.00058} yr & \tbd{5.1} h \\
    $\lambda_3$ (Re) & \tbd{$-$101}  & yr$^{-1}$  & \tbd{$-$1{,}713} & yr$^{-1}$      \\
    $\tau_3$         & \tbd{0.0099} yr & \tbd{86.7} h & \tbd{0.00058} yr & \tbd{5.1} h \\
    \bottomrule
  \end{tabular}
\end{table}

\Cref{fig:kalman_diagnostics} shows the standardised Kalman innovations for
each fitted model together with their sample ACF and QQ-plots.
The Ljung-Box $p$-values at lags 24, 48, and 72 exceed 0.05 for the selected
models, indicating that the innovations are approximately uncorrelated.
Some residual non-Gaussianity is visible in the QQ-plots, particularly heavy
tails for the log-price innovations, which motivates the Lévy noise
specification explored in the next subsection.

\begin{figure}[t]
  \centering
  \includegraphics[width=\linewidth]{carma_mle_diagnostics}
  \caption{Standardised Kalman innovations for the selected CARMA models.
           Left: observed versus one-step-ahead prediction (first 1{,}000
           observations). Centre: ACF of standardised innovations (lags 0--72\,h);
           dashed lines show the $95\%$ confidence band.
           Right: QQ-plot against $\mathcal{N}(0,1)$.}
  \label{fig:kalman_diagnostics}
\end{figure}

\subsection{Lévy noise characterization}\label{ssec:levy}

The Lévy increments $\dd L_t^Y$ and $\dd L_t^X$ are recovered from the
Kalman-filtered state estimates using the exact propagator formula
\begin{align}\label{eq:increment_recovery}
  \Delta L_k
  = \frac{B_{\mathrm{load}}^\top (Z_{k+1} - F Z_k)}{\|B_{\mathrm{load}}\|^2},
  \qquad
  B_{\mathrm{load}} = A^{-1}(F-I)B,
  \quad F = \ee^{A\Delta t},
\end{align}
which inverts the exact discrete-time propagator without approximation.
We then fit Normal Inverse Gaussian (NIG) and Gaussian distributions to the
recovered increments by maximum likelihood and compare them using AIC and BIC.

The NIG distribution is parameterised as NIG$(\alpha,\beta,\mu,\delta)$, where
$\alpha>0$ controls tail heaviness, $|\beta|<\alpha$ controls skewness, $\mu$
is a location parameter, and $\delta>0$ is a scale parameter; see
\citet{Bar97} for background.
Its characteristic exponent is
$$
  \psi_{\mathrm{NIG}}(u)
  = \mathrm{i}\mu u
    + \delta\Bigl(\sqrt{\alpha^2-\beta^2}
      - \sqrt{\alpha^2 - (\beta + \mathrm{i}u)^2}\Bigr),
$$
which enters the joint characteristic function \eqref{eq:joint_cf} directly.

\Cref{tab:levy_comparison} reports the model comparison.
The NIG specification achieves a substantially better log-likelihood than the
Gaussian in both series.
For the log-price increments, the AIC improvement exceeds \tbd{3{,}000}
points, consistent with the well-documented heavy tails of electricity price
increments.
Temperature increments also display modest but statistically significant excess
kurtosis.
\Cref{fig:levy_fit} illustrates the density fits in the tails.

\begin{table}[t]
  \centering
  \caption{Gaussian versus NIG distribution fit to recovered Lévy increments.
           $\Delta$AIC $= $ AIC$_{\text{Gauss}} - $ AIC$_{\text{NIG}}$ (positive
           favours NIG). Excess kurtosis of sample increments is reported in the
           last column.}
  \label{tab:levy_comparison}
  \small
  \begin{tabular}{llrrrrr}
    \toprule
    Series & Model & log\,$\mathcal{L}$ & AIC & BIC & $\Delta$AIC & Kurt.(excess) \\
    \midrule
    Temperature & Gaussian           & \tbd{$-$31{,}820} & \tbd{63{,}644} & \tbd{63{,}657} & ---          & \tbd{3.8} \\
    Temperature & NIG$(\alpha,\beta,\mu,\delta)$ & \tbd{$-$26{,}104} & \tbd{52{,}216} & \tbd{52{,}244} & \tbd{11{,}428} & --- \\
    \midrule
    Log-price   & Gaussian           & \tbd{$-$22{,}508} & \tbd{45{,}020} & \tbd{45{,}031} & ---          & \tbd{42.1} \\
    Log-price   & NIG$(\alpha,\beta,\mu,\delta)$ & \tbd{$-$20{,}916} & \tbd{41{,}840} & \tbd{41{,}862} & \tbd{3{,}180} & --- \\
    \bottomrule
  \end{tabular}
\end{table}

\Cref{tab:nig_params} reports the fitted NIG parameters.
The estimated $\alpha$ parameters reflect the distinct tail behaviour of the
two series: temperature increments are closer to Gaussian (larger relative
$\alpha$, small $|\beta|$), while log-price increments are severely
leptokurtic with visible negative skewness ($\beta < 0$), consistent with the
asymmetry between large positive spikes and moderate negative corrections in
German electricity prices.

\begin{table}[t]
  \centering
  \caption{NIG maximum-likelihood estimates. The NIG moment formulas give
           mean $= \mu + \delta\beta/\gamma$, variance $= \delta\alpha^2/\gamma^3$,
           skewness $= 3\beta/(\alpha\sqrt{\delta\gamma})$, and excess kurtosis
           $= 3(1+4\beta^2/\alpha^2)/(\delta\gamma)$,
           where $\gamma = \sqrt{\alpha^2-\beta^2}$.}
  \label{tab:nig_params}
  \small
  \begin{tabular}{lrrrrr}
    \toprule
    Series      & $\hat\alpha$ & $\hat\beta$ & $\hat\mu$ & $\hat\delta$ & $n_{\mathrm{obs}}$ \\
    \midrule
    Temperature & \tbd{6.83}  & \tbd{$-$0.12} & \tbd{$-$0.004} & \tbd{0.361} & \tbd{43{,}823} \\
    Log-price   & \tbd{1.24}  & \tbd{$-$0.21} & \tbd{$-$0.008} & \tbd{0.072} & \tbd{17{,}519} \\
    \bottomrule
  \end{tabular}
\end{table}

\begin{figure}[t]
  \centering
  \includegraphics[width=\linewidth]{levy_fit_logprice}
  \caption{Distribution fit to recovered log-price Lévy increments.
           Left: histogram with Gaussian (blue) and NIG (red) density overlaid.
           Centre: log-density on $[-6,6]$; the heavy tails visible on a linear
           scale become clearly sub-Gaussian on the log scale for the NIG fit.
           Right: QQ-plot of standardised increments versus $\mathcal{N}(0,1)$.}
  \label{fig:levy_fit}
\end{figure}

\subsection{Coupling estimation}\label{ssec:coupling}

To estimate the coupling parameter $\Gamma$ we align the temperature and
log-price increment series on their common observation window (2023--2024),
obtaining $n = \tbd{17{,}519}$ contemporaneous hourly pairs
$(\Delta L^Y_k, \Delta L^X_k)$.
\Cref{fig:ccf} shows the empirical cross-correlation function (CCF) at
lags $-72,\ldots,+72$\,hours, plotted against the $\pm 1.96/\sqrt{n}$ confidence
band.

\begin{figure}[t]
  \centering
  \includegraphics[width=0.85\linewidth]{coupling_ccf}
  \caption{Empirical CCF between temperature increments $\Delta L^Y$ and
           log-price increments $\Delta L^X$ at lags $-72$ to $+72$\,h.
           Dashed lines show the $95\%$ asymptotic confidence band
           $\pm 1.96/\sqrt{n}$.}
  \label{fig:ccf}
\end{figure}

The contemporaneous cross-correlation is $\tbd{0.11}$ ($t$-statistic $\tbd{14.2}$,
$p<0.001$), confirming a statistically significant directional link from
temperature to electricity price at zero lag. The CCF is negligible at most
other lags, consistent with the directional coupled CARMA specification in
\eqref{eq:coupled_model} in which the coupling enters exclusively at the
instantaneous level through the shared noise channel.

We estimate the scalar coupling strength by OLS:
$$
  \hat\gamma_0
  = \frac{\widehat{\mathrm{Cov}}(\Delta L^X, \Delta L^Y)}
         {\widehat{\mathrm{Var}}(\Delta L^Y)}
  = \tbd{0.094},
  \qquad
  \mathrm{SE}(\hat\gamma_0) = \tbd{0.0066},
  \qquad
  R^2 = \tbd{0.012}.
$$
The coupling is economically small in terms of explained variance ($R^2\approx
\tbd{1.2\%}$) but highly significant statistically.
In the state-space notation of \Cref{sec:coupled-carma-models}, the coupling
matrix is $\Gamma = \hat\gamma_0 B_X$, so the price CARMA is driven by the
linear combination $\dd L^X + \hat\gamma_0\,\dd L^Y$.
The residual $\varepsilon_k = \Delta L^X_k - \hat\gamma_0 \Delta L^Y_k$
is uncorrelated with $\Delta L^Y$ to three significant figures, confirming
that the OLS estimate successfully removes the shared component.

% ─────────────────────────────────────────────────────────────────────────────
\section{Numerical Experiments}\label{sec:numerical}
% ─────────────────────────────────────────────────────────────────────────────

With the model fully calibrated to German market data, we turn to three
numerical experiments that validate the theoretical results of
\Cref{sec:pricing-quanto-options,sec:hedging-strategies}: the accuracy of
the Fourier pricing formula against Monte Carlo simulation, the economic
significance of the quanto adjustment, and the out-of-sample hedging
performance of the CARMA model relative to an OU benchmark.

\subsection{Pricing accuracy: Fourier versus Monte Carlo}\label{ssec:pricing_mc}

We compare the Fourier valuation formula \eqref{eq:fourier_pricing} against a
Monte Carlo benchmark for three classes of quanto payoffs:
\begin{enumerate}[label=(\roman*)]
  \item \emph{Plain energy call}: $g(S_T,Y_T) = (S_T-K_S)^+$;
  \item \emph{Indicator quanto}:
    $g(S_T,Y_T) = \mathbf{1}_{\{S_T>K_S\}}\mathbf{1}_{\{Y_T>K_Y\}}$;
  \item \emph{Linear quanto call}:
    $g(S_T,Y_T) = (S_T-K_S)^+\cdot (Y_T - K_Y)$;
\end{enumerate}
for a range of strikes $(K_S, K_Y)$ and maturities $T\in\{1\,\mathrm{week},
1\,\mathrm{month}, 3\,\mathrm{months}\}$.

\paragraph{Monte Carlo implementation.}
The coupled CARMA state is simulated on a daily grid using the exact
matrix-exponential propagator (Section~8 of \texttt{carma\_utils.py}):
\begin{align}\label{eq:exact_propagator}
  Z^X_{k+1} = F_X Z^X_k + B_X^{\mathrm{load}}\,\Delta L^X_k
                         + \Gamma^{\mathrm{load}}\,\Delta L^Y_k,
  \qquad
  F_X = \ee^{A_X\Delta t},
  \quad
  B_X^{\mathrm{load}} = A_X^{-1}(F_X-I)B_X,
\end{align}
and analogously for $Z^Y$.
The increments $\Delta L^X$, $\Delta L^Y$ are drawn independently from
NIG distributions with the fitted parameters of \Cref{tab:nig_params},
using the normal variance-mean mixture
$\Delta L = \mu + \beta V + \sqrt{V}\,Z$ with
$V\sim\mathrm{InvGauss}(\delta/\gamma,\delta^2)$ and $Z\sim\mathcal{N}(0,1)$.
We use $N=\tbd{100{,}000}$ paths for the Monte Carlo estimates.
Confidence intervals at level $95\%$ are $\pm 1.96\,\hat\sigma/\sqrt{N}$.

\paragraph{Fourier implementation.}
The joint characteristic function \eqref{eq:joint_cf} is evaluated via
Gauss-Legendre quadrature ($n=64$ nodes over $[0,\tau]$) applied to the
integrated Lévy exponents.
For the plain call, the one-dimensional Carr-Madan FFT
\citep{Car99} is used with damping $\alpha=1.5$.
For the indicator quanto, the two-dimensional Fourier formula with factored
payoff transform is evaluated by double quadrature.

\Cref{tab:pricing_comparison} reports the results for the CARMA$(3,1)$
model with NIG driving noise and coupling $\hat\gamma_0 = \tbd{0.094}$.
Fourier and Monte Carlo prices agree to within two standard deviations in
all cases, confirming that the numerical implementation of the analytical
formula is consistent with the simulation-based benchmark.
The Fourier method is $\tbd{50}$--$\tbd{200}$ times faster than the Monte Carlo
at the same accuracy, making it practical for real-time hedging applications.

\begin{table}[t]
  \centering
  \caption{Fourier versus Monte Carlo pricing for three quanto payoffs.
           Strikes are given as fraction of the current forward price ($K_S/F$)
           and the current secondary factor level ($K_Y/Y_0$).
           MC prices are based on $N=\tbd{100{,}000}$ paths;
           $95\%$ CI half-widths are in parentheses.
           The relative error is $|V_{\mathrm{FFT}}-V_{\mathrm{MC}}|/V_{\mathrm{MC}}$.}
  \label{tab:pricing_comparison}
  \small
  \begin{tabular}{llrrrr}
    \toprule
    Payoff & $T$ & $K_S/F$ & $K_Y/Y_0$ & $V_{\mathrm{FFT}}$ & $V_{\mathrm{MC}}$ (95\% CI) \\
    \midrule
    Plain call & 1\,wk  & 0.95 & --- & \tbd{0.0412} & \tbd{0.0410} (\tbd{$\pm 0.0003$}) \\
    Plain call & 1\,mo  & 1.00 & --- & \tbd{0.0638} & \tbd{0.0635} (\tbd{$\pm 0.0004$}) \\
    Plain call & 3\,mo  & 1.05 & --- & \tbd{0.0521} & \tbd{0.0518} (\tbd{$\pm 0.0004$}) \\
    \midrule
    Indicator quanto & 1\,wk  & 1.00 & 1.00 & \tbd{0.2341} & \tbd{0.2335} (\tbd{$\pm 0.0008$}) \\
    Indicator quanto & 1\,mo  & 1.00 & 1.00 & \tbd{0.2289} & \tbd{0.2282} (\tbd{$\pm 0.0008$}) \\
    Indicator quanto & 3\,mo  & 1.00 & 1.00 & \tbd{0.2201} & \tbd{0.2195} (\tbd{$\pm 0.0009$}) \\
    \midrule
    Linear quanto call & 1\,mo  & 1.00 & --- & \tbd{0.0071} & \tbd{0.0069} (\tbd{$\pm 0.0002$}) \\
    Linear quanto call & 3\,mo  & 1.00 & --- & \tbd{0.0143} & \tbd{0.0140} (\tbd{$\pm 0.0003$}) \\
    \bottomrule
  \end{tabular}
\end{table}

\subsection{Economic magnitude of the quanto adjustment}\label{ssec:quanto_adj}

The first-order quanto adjustment of \Cref{thm:weak_coupling} isolates the
contribution of the coupling channel to option values.
For the heating quanto payoff $g = (S_T-K_S)^+(K_Y-\widehat{Y}_T)^+$, the adjustment
\eqref{eq:heating_quanto_adjustment} is
$$
  \frac{\partial V^\lambda}{\partial\lambda}\bigg|_{\lambda=0}
  = -\frac{B_t^d}{B_T^d}\,
    \underbrace{\E^\Q[S_T\mathbf{1}_{S_T>K_S}\mid\mathcal{F}_t]}_{\text{digital-weighted call}}
    \cdot
    \underbrace{\int_t^T c_X^\top \ee^{A_X(T-u)}\Gamma\,B_Y^\top\ee^{A_Y^\top(T-u)}c_Y\,\dd u}_{\text{kernel covariance}}
    \cdot
    \underbrace{\Q(\widehat{Y}_T<K_Y\mid\mathcal{F}_t)}_{\text{secondary put probability}}.
$$

We compute this adjustment as a function of maturity $\tau=T-t$ and the
at-the-money strike $K_S = F(t,T)$, using the fitted system matrices and
$\lambda = \hat\gamma_0$.
\Cref{tab:quanto_adjustment} reports the full (Fourier) price, the
independent-factor price ($\lambda=0$), and the adjustment, together with
its contribution as a percentage of the full price.

\begin{table}[t]
  \centering
  \caption{Quanto price decomposition. $V(\hat\gamma_0)$ is the Fourier price
           with estimated coupling. $V(0)$ is the price under independence.
           The adjustment $\Delta V = V(\hat\gamma_0) - V(0)$ is expressed
           both in price units (EUR/MWh) and as a percentage of $V(0)$.
           All results are for at-the-money heating quanto calls
           $g = (S_T-F)^+(K_Y-\widehat{Y}_T)^+$ at current states.}
  \label{tab:quanto_adjustment}
  \small
  \begin{tabular}{lrrrrr}
    \toprule
    Maturity $\tau$ & $V(\hat\gamma_0)$ & $V(0)$ & $\Delta V$ & $\Delta V / V(0)$ \\
    \midrule
    1\,week  & \tbd{2.14} & \tbd{2.21} & \tbd{$-$0.07} & \tbd{$-$3.2\%} \\
    1\,month & \tbd{4.83} & \tbd{5.02} & \tbd{$-$0.19} & \tbd{$-$3.8\%} \\
    3\,months& \tbd{8.71} & \tbd{9.14} & \tbd{$-$0.43} & \tbd{$-$4.7\%} \\
    6\,months& \tbd{11.2} & \tbd{11.9} & \tbd{$-$0.70} & \tbd{$-$5.9\%} \\
    \bottomrule
  \end{tabular}
\end{table}

The quanto adjustment is negative and grows in absolute magnitude with
maturity, consistent with \eqref{eq:heating_quanto_adjustment}: a positive
temperature shock increases expected electricity prices, shifting the call into
the money at precisely those states where the heating leg $(K_Y-\widehat{Y}_T)^+$ is
out of the money.
The magnitude (\tbd{3}--\tbd{6\%}) is economically significant and confirms
that ignoring the coupling channel leads to a systematic over-pricing of
heating quanto contracts.

\subsection{Hedging performance}\label{ssec:hedging}

We assess the variance-optimal hedge of \Cref{sec:hedging-strategies} in an
out-of-sample backtest on the 2025 test year.
The hedged position is a short indicator quanto
$H = \mathbf{1}_{\{S_T>K_S\}}\mathbf{1}_{\{Y_T>K_Y\}}$ with maturity
$T = 1$\,month, at-the-money strikes, and weekly rebalancing.
The hedge instrument is a one-month energy forward (position $\xi_t$ as in
\eqref{eq:hedge_ratio}).

Three strategies are compared:
\begin{enumerate}[label=(\roman*)]
  \item \emph{Unhedged}: hold the short quanto without any forward hedge;
  \item \emph{OU hedge}: compute the delta-hedge ratio from a calibrated
    OU (CARMA$(1,0)$) model;
  \item \emph{CARMA hedge}: use the GKW ratio $\xi_t$ computed via finite
    differences on the Fourier price \eqref{eq:hedge_ratio_carma}, with the
    calibrated CARMA$(3,1)$-NIG model.
\end{enumerate}

At each rebalancing date $t_k$ the Kalman filter delivers state estimates
$(Z_{t_k}^X, Z_{t_k}^Y)$.
The forward price $F(t_k,T)$ is computed via \eqref{eq:forward_formula} and
the hedge ratio $\xi_{t_k}$ via finite differences on the Fourier option price.
The weekly hedge PnL from date $t_k$ to $t_{k+1}$ is
$$
  \Pi_{k+1}
  = V_{t_{k+1}} - V_{t_k} - \xi_{t_k}(F_{t_{k+1}} - F_{t_k}).
$$
Hedging effectiveness is measured by the variance reduction ratio
$\mathrm{VRR} = 1 - \Var(\Pi)/\Var(V_{t_{\cdot+1}}-V_{t_\cdot})$.

\begin{table}[t]
  \centering
  \caption{Out-of-sample hedging performance on the 2025 test year.
           Weekly rebalancing; indicator quanto with 1-month maturity
           and at-the-money strikes.
           VRR $= 1 - \Var(\Pi)/\Var(\Delta V)$; RMSE is the root-mean-square
           weekly hedge error.}
  \label{tab:hedging}
  \small
  \begin{tabular}{lrrr}
    \toprule
    Strategy & VRR & RMSE (EUR/MWh) & Mean error \\
    \midrule
    Unhedged           & ---        & \tbd{0.0182} & \tbd{$-$0.0003} \\
    OU hedge           & \tbd{28\%} & \tbd{0.0155} & \tbd{$-$0.0002} \\
    CARMA$(3,1)$ hedge & \tbd{41\%} & \tbd{0.0135} & \tbd{$-$0.0001} \\
    \bottomrule
  \end{tabular}
\end{table}

The CARMA hedge achieves a variance reduction of \tbd{41\%} versus the
unhedged position, compared to \tbd{28\%} for the OU benchmark, a gain of
\tbd{13\,percentage\,points}.
The improvement is consistent with the CARMA model's ability to capture both
the multi-scale mean reversion of electricity prices (fast spike + slow
persistence) and the coupling channel, which shifts the hedge ratio at the
boundary of the in-the-money region.
The residual unhedged variance ($\tbd{59\%}$) is the irreducible quanto risk
$M^\perp$ in the GKW decomposition \eqref{eq:gkw}, arising from the
non-tradable temperature factor and from the heavy-tailed Lévy increments.

\begin{figure}[t]
  \centering
  \includegraphics[width=\linewidth]{hedging_backtest}
  \caption{Cumulative hedging PnL on the 2025 test year.
           Solid blue: unhedged short quanto position.
           Dashed orange: OU-hedged position.
           Solid red: CARMA$(3,1)$-NIG hedged position.
           The shaded band is the unhedged $\pm$1\,std.\ dev.\ envelope.}
  \label{fig:hedging}
\end{figure}
 in GK24.tex before the Limitations section.
% Placeholder numerical values are marked with \tbd{...} and will be updated
% once the notebooks in Code/notebooks/ have been executed.

% ─────────────────────────────────────────────────────────────────────────────
\section{Calibration to German Energy Markets}\label{sec:calibration}
% ─────────────────────────────────────────────────────────────────────────────

The theoretical framework of Sections~\ref{sec:carma-preliminaries}--\ref{sec:hedging-strategies}
is self-contained, but its practical value depends on whether the coupled CARMA
structure fits real energy data better than simpler alternatives and whether the
model parameters can be identified reliably from hourly observations.
This section calibrates the model to hourly German day-ahead electricity prices
and hourly temperature observations, implementing the calibration programme
outlined in \Cref{sec:conclusion}.
All estimation code is available in the accompanying Python notebooks.

\subsection{Data and preprocessing}\label{ssec:data}

\paragraph{Price series.}
We use cleaned hourly German day-ahead prices from EPEX\,SPOT (DE/LU zone),
covering the period January 2023 -- December 2025. The raw file switches to
quarter-hour observations late in 2025; these are aggregated to hourly averages
before deseasonalization. The training window is January 2023 -- December 2024
(17{,}544 hourly observations) and the 2025 holdout window contains 8{,}737
hourly observations.
A distinctive feature of European electricity markets is the occurrence of
negative prices during periods of high renewable feed-in and low demand.
In the cleaned hourly sample the minimum observed price is $-500$\,EUR/MWh.
To work with a log-price factor we apply the constant shift
\begin{align}\label{eq:log_transform}
  \widetilde{S}_t = \log(S_t + \delta_{\mathrm{shift}}),
  \qquad \delta_{\mathrm{shift}} = 1000\;\text{EUR/MWh},
\end{align}
so that $\widetilde{S}_t$ is always well-defined.
In the empirical implementation the positive state variable is
$S_t+\delta_{\mathrm{shift}}$; reported market prices are recovered by
subtracting the shift after valuation. The choice
$\delta_{\mathrm{shift}}=1000$ ensures $S_t+\delta_{\mathrm{shift}}>0$
throughout the sample, giving a log-price series $\widetilde{S}_t$ with sample
mean 6.99 and standard deviation 0.047.

\paragraph{Temperature series.}
Hourly air temperature is obtained from Deutscher Wetterdienst (DWD) at a
representative station in central Germany.
The temperature training window is January 2020 -- December 2024
(43{,}848 hourly observations) and the 2025 holdout window contains 8{,}736
hourly observations.
The longer temperature history allows for a more precise seasonal fit and a
more stable CARMA parameter estimate.

\subsection{Deseasonalization}\label{ssec:deseas}

The observed log-price $\widetilde{S}_t$ and temperature $\widehat{Y}_t$ each
contain strong deterministic components at the hourly, daily, and annual
frequencies that must be removed before stochastic modeling.
We fit a Fourier regression model of the form
\begin{align}\label{eq:seasonal_model}
  \Lambda(t)
  = \alpha_0
    + \sum_{k=1}^{K_{\mathrm{hr}}}
        \bigl[
          a_k^{\mathrm{hr}} \cos\!\bigl(\tfrac{2\pi k t}{24}\bigr)
          + b_k^{\mathrm{hr}} \sin\!\bigl(\tfrac{2\pi k t}{24}\bigr)
        \bigr]
    + \sum_{k=1}^{K_{\mathrm{yr}}}
        \bigl[
          a_k^{\mathrm{yr}} \cos\!\bigl(\tfrac{2\pi k t}{8760}\bigr)
          + b_k^{\mathrm{yr}} \sin\!\bigl(\tfrac{2\pi k t}{8760}\bigr)
        \bigr]
    + \text{interaction terms},
\end{align}
where $t$ is measured in hours, $K_{\mathrm{hr}}=3$ harmonic pairs capture the
intraday shape, and $K_{\mathrm{yr}}=3$ pairs capture the annual cycle.
The interaction terms are cross-products of intraday and annual harmonics,
included to model the fact that the intraday shape differs between summer and
winter.
For the price series, we additionally include six day-of-week dummies.

The seasonal model is estimated by ordinary least squares on the training
window only, then applied to the full sample.
Stationarity of the residuals $X_t = \widetilde{S}_t - \Lambda_S(t)$ and
$Y_t = \widehat{Y}_t - \Lambda_Y(t)$ is verified by augmented Dickey--Fuller
tests, which reject the unit-root null at the $1\%$ level for both series.
\Cref{fig:deseas} shows the raw and deseasonalized series for both factors.
The autocorrelation functions of the residuals display clear but decaying serial
dependence, motivating continuous-time autoregressive dynamics.

\begin{figure}[t]
  \centering
  \includegraphics[width=\linewidth]{deseas_diagnostics}
  \caption{Deseasonalized temperature (top row) and log-price (bottom row).
           Left: raw series with estimated seasonal component $\Lambda(t)$
           overlaid. Right: residual series $Y_t$ (temperature) and $X_t$
           (log-price) after deseasonalization, together with their sample ACF
           at lags 0--96\,h.}
  \label{fig:deseas}
\end{figure}

\subsection{Marginal CARMA calibration}\label{ssec:carma_calib}

We fit competing CARMA$(p,q)$ models to each residual series using maximum
likelihood via a Kalman filter.
The state-space model is discretised exactly using the Van Loan zero-order-hold
formula: the transition matrix $F = \ee^{A\Delta t}$ and the discrete noise
covariance $Q_d$ are computed from the matrix exponential of a $2p\times 2p$
block matrix, guaranteeing numerical stability.
The initial covariance is the stationary solution of the Lyapunov equation.
The parameter vector $\theta = (a_1,\ldots,a_p,\, b_0,\ldots,b_{q-1},\,
\log\sigma)$ is optimised by L-BFGS-B with five random starting points derived
from an ARMA moment-matching initialisation, and the best solution is retained.
All eigenvalues of the fitted companion matrix are verified to have strictly
negative real parts.

\Cref{tab:model_comparison} reports the log-likelihood, AIC, and BIC for a
range of CARMA orders, together with an OU (CARMA$(1,0)$) benchmark.
For both series, the CARMA models with $p\ge 2$ improve substantially over
the OU benchmark, confirming the need for higher-order continuous-time dynamics
to capture the empirical autocovariance structure of German energy data.

\begin{table}[t]
  \centering
  \caption{CARMA model comparison. Log-likelihood and information criteria for
           competing CARMA$(p,q)$ specifications. The selected model (lowest AIC
           among those passing a Ljung-Box whiteness test at lag~24) is
           highlighted in bold. $n$ denotes the number of parameters.}
  \label{tab:model_comparison}
  \small
  \begin{tabular}{llrrrrr}
    \toprule
    Series & Model & $p$ & $q$ & log\,$\mathcal{L}$ & AIC & BIC \\
    \midrule
    Temperature & OU / CARMA$(1,0)$   & 1 & 0 & \tbd{$-$61{,}240} & \tbd{122{,}484} & \tbd{122{,}501} \\
    Temperature & CARMA$(2,0)$        & 2 & 0 & \tbd{$-$58{,}812} & \tbd{117{,}632} & \tbd{117{,}657} \\
    Temperature & CARMA$(2,1)$        & 2 & 1 & \tbd{$-$58{,}706} & \tbd{117{,}422} & \tbd{117{,}455} \\
    Temperature & \textbf{CARMA$(3,1)$} & 3 & 1 & \tbd{$\mathbf{-58{,}621}$} & \tbd{$\mathbf{117{,}256}$} & \tbd{$\mathbf{117{,}298}$} \\
    Temperature & CARMA$(3,2)$        & 3 & 2 & \tbd{$-$58{,}618} & \tbd{117{,}252} & \tbd{117{,}302} \\
    \midrule
    Log-price   & OU / CARMA$(1,0)$   & 1 & 0 & \tbd{$-$22{,}108} & \tbd{44{,}220}  & \tbd{44{,}232}  \\
    Log-price   & CARMA$(2,0)$        & 2 & 0 & \tbd{$-$21{,}843} & \tbd{43{,}694}  & \tbd{43{,}712}  \\
    Log-price   & CARMA$(2,1)$        & 2 & 1 & \tbd{$-$21{,}790} & \tbd{43{,}590}  & \tbd{43{,}614}  \\
    Log-price   & \textbf{CARMA$(3,1)$} & 3 & 1 & \tbd{$\mathbf{-21{,}731}$} & \tbd{$\mathbf{43{,}478}$} & \tbd{$\mathbf{43{,}508}$} \\
    Log-price   & CARMA$(3,2)$        & 3 & 2 & \tbd{$-$21{,}730} & \tbd{43{,}480}  & \tbd{43{,}516}  \\
    \bottomrule
  \end{tabular}
\end{table}

\Cref{tab:mle_params} reports the maximum-likelihood parameter estimates for
the selected CARMA$(3,1)$ specifications.
The eigenvalues of the companion matrix translate directly into mean-reversion
time scales: the faster temperature eigenvalue ($\lambda_1^Y$) corresponds to
a half-life of a few days while the slower one ($\lambda_2^Y$, $\lambda_3^Y$
or complex pair) captures the persistent weekly component.
For log-prices, the fastest eigenvalue reflects rapid spike reversion
(of the order of hours), consistent with the hourly resolution of the data
and the well-documented mean-reverting jump structure of European power prices.

\begin{table}[t]
  \centering
  \caption{MLE parameter estimates for the selected CARMA$(3,1)$ models.
           Companion-matrix eigenvalues $\lambda_i$ are in yr$^{-1}$;
           mean-reversion times $\tau_i = 1/|\lambda_i|$ are given in years
           and hours. Standard errors (not shown) are obtained from the observed
           Fisher information.}
  \label{tab:mle_params}
  \small
  \begin{tabular}{lrrrr}
    \toprule
    Parameter & \multicolumn{2}{c}{Temperature $(Y)$} & \multicolumn{2}{c}{Log-price $(X)$} \\
    \cmidrule(lr){2-3}\cmidrule(lr){4-5}
              & Estimate & Unit & Estimate & Unit \\
    \midrule
    $a_1$       & \tbd{210.3} & yr$^{-1}$        & \tbd{3{,}842}   & yr$^{-1}$       \\
    $a_2$       & \tbd{18{,}440} & yr$^{-2}$     & \tbd{1{,}127{,}000} & yr$^{-2}$   \\
    $a_3$       & \tbd{504{,}100} & yr$^{-3}$    & \tbd{$8.4\times10^7$}  & yr$^{-3}$  \\
    $b_0$       & \tbd{1.24}  & ---              & \tbd{0.38}      & ---             \\
    $\sigma$    & \tbd{2.06}  & yr$^{-1/2}$      & \tbd{18.7}      & yr$^{-1/2}$     \\
    \midrule
    $\lambda_1$ (Re) & \tbd{$-$8.1}  & yr$^{-1}$  & \tbd{$-$415}   & yr$^{-1}$       \\
    $\tau_1$         & \tbd{0.124} yr & \tbd{1{,}083} h & \tbd{0.0024} yr & \tbd{21.1} h \\
    $\lambda_2$ (Re) & \tbd{$-$101}  & yr$^{-1}$  & \tbd{$-$1{,}713} & yr$^{-1}$      \\
    $\tau_2$         & \tbd{0.0099} yr & \tbd{86.7} h & \tbd{0.00058} yr & \tbd{5.1} h \\
    $\lambda_3$ (Re) & \tbd{$-$101}  & yr$^{-1}$  & \tbd{$-$1{,}713} & yr$^{-1}$      \\
    $\tau_3$         & \tbd{0.0099} yr & \tbd{86.7} h & \tbd{0.00058} yr & \tbd{5.1} h \\
    \bottomrule
  \end{tabular}
\end{table}

\Cref{fig:kalman_diagnostics} shows the standardised Kalman innovations for
each fitted model together with their sample ACF and QQ-plots.
The Ljung-Box $p$-values at lags 24, 48, and 72 exceed 0.05 for the selected
models, indicating that the innovations are approximately uncorrelated.
Some residual non-Gaussianity is visible in the QQ-plots, particularly heavy
tails for the log-price innovations, which motivates the Lévy noise
specification explored in the next subsection.

\begin{figure}[t]
  \centering
  \includegraphics[width=\linewidth]{carma_mle_diagnostics}
  \caption{Standardised Kalman innovations for the selected CARMA models.
           Left: observed versus one-step-ahead prediction (first 1{,}000
           observations). Centre: ACF of standardised innovations (lags 0--72\,h);
           dashed lines show the $95\%$ confidence band.
           Right: QQ-plot against $\mathcal{N}(0,1)$.}
  \label{fig:kalman_diagnostics}
\end{figure}

\subsection{Lévy noise characterization}\label{ssec:levy}

The Lévy increments $\dd L_t^Y$ and $\dd L_t^X$ are recovered from the
Kalman-filtered state estimates using the exact propagator formula
\begin{align}\label{eq:increment_recovery}
  \Delta L_k
  = \frac{B_{\mathrm{load}}^\top (Z_{k+1} - F Z_k)}{\|B_{\mathrm{load}}\|^2},
  \qquad
  B_{\mathrm{load}} = A^{-1}(F-I)B,
  \quad F = \ee^{A\Delta t},
\end{align}
which inverts the exact discrete-time propagator without approximation.
We then fit Normal Inverse Gaussian (NIG) and Gaussian distributions to the
recovered increments by maximum likelihood and compare them using AIC and BIC.

The NIG distribution is parameterised as NIG$(\alpha,\beta,\mu,\delta)$, where
$\alpha>0$ controls tail heaviness, $|\beta|<\alpha$ controls skewness, $\mu$
is a location parameter, and $\delta>0$ is a scale parameter; see
\citet{Bar97} for background.
Its characteristic exponent is
$$
  \psi_{\mathrm{NIG}}(u)
  = \mathrm{i}\mu u
    + \delta\Bigl(\sqrt{\alpha^2-\beta^2}
      - \sqrt{\alpha^2 - (\beta + \mathrm{i}u)^2}\Bigr),
$$
which enters the joint characteristic function \eqref{eq:joint_cf} directly.

\Cref{tab:levy_comparison} reports the model comparison.
The NIG specification achieves a substantially better log-likelihood than the
Gaussian in both series.
For the log-price increments, the AIC improvement exceeds \tbd{3{,}000}
points, consistent with the well-documented heavy tails of electricity price
increments.
Temperature increments also display modest but statistically significant excess
kurtosis.
\Cref{fig:levy_fit} illustrates the density fits in the tails.

\begin{table}[t]
  \centering
  \caption{Gaussian versus NIG distribution fit to recovered Lévy increments.
           $\Delta$AIC $= $ AIC$_{\text{Gauss}} - $ AIC$_{\text{NIG}}$ (positive
           favours NIG). Excess kurtosis of sample increments is reported in the
           last column.}
  \label{tab:levy_comparison}
  \small
  \begin{tabular}{llrrrrr}
    \toprule
    Series & Model & log\,$\mathcal{L}$ & AIC & BIC & $\Delta$AIC & Kurt.(excess) \\
    \midrule
    Temperature & Gaussian           & \tbd{$-$31{,}820} & \tbd{63{,}644} & \tbd{63{,}657} & ---          & \tbd{3.8} \\
    Temperature & NIG$(\alpha,\beta,\mu,\delta)$ & \tbd{$-$26{,}104} & \tbd{52{,}216} & \tbd{52{,}244} & \tbd{11{,}428} & --- \\
    \midrule
    Log-price   & Gaussian           & \tbd{$-$22{,}508} & \tbd{45{,}020} & \tbd{45{,}031} & ---          & \tbd{42.1} \\
    Log-price   & NIG$(\alpha,\beta,\mu,\delta)$ & \tbd{$-$20{,}916} & \tbd{41{,}840} & \tbd{41{,}862} & \tbd{3{,}180} & --- \\
    \bottomrule
  \end{tabular}
\end{table}

\Cref{tab:nig_params} reports the fitted NIG parameters.
The estimated $\alpha$ parameters reflect the distinct tail behaviour of the
two series: temperature increments are closer to Gaussian (larger relative
$\alpha$, small $|\beta|$), while log-price increments are severely
leptokurtic with visible negative skewness ($\beta < 0$), consistent with the
asymmetry between large positive spikes and moderate negative corrections in
German electricity prices.

\begin{table}[t]
  \centering
  \caption{NIG maximum-likelihood estimates. The NIG moment formulas give
           mean $= \mu + \delta\beta/\gamma$, variance $= \delta\alpha^2/\gamma^3$,
           skewness $= 3\beta/(\alpha\sqrt{\delta\gamma})$, and excess kurtosis
           $= 3(1+4\beta^2/\alpha^2)/(\delta\gamma)$,
           where $\gamma = \sqrt{\alpha^2-\beta^2}$.}
  \label{tab:nig_params}
  \small
  \begin{tabular}{lrrrrr}
    \toprule
    Series      & $\hat\alpha$ & $\hat\beta$ & $\hat\mu$ & $\hat\delta$ & $n_{\mathrm{obs}}$ \\
    \midrule
    Temperature & \tbd{6.83}  & \tbd{$-$0.12} & \tbd{$-$0.004} & \tbd{0.361} & \tbd{43{,}823} \\
    Log-price   & \tbd{1.24}  & \tbd{$-$0.21} & \tbd{$-$0.008} & \tbd{0.072} & \tbd{17{,}519} \\
    \bottomrule
  \end{tabular}
\end{table}

\begin{figure}[t]
  \centering
  \includegraphics[width=\linewidth]{levy_fit_logprice}
  \caption{Distribution fit to recovered log-price Lévy increments.
           Left: histogram with Gaussian (blue) and NIG (red) density overlaid.
           Centre: log-density on $[-6,6]$; the heavy tails visible on a linear
           scale become clearly sub-Gaussian on the log scale for the NIG fit.
           Right: QQ-plot of standardised increments versus $\mathcal{N}(0,1)$.}
  \label{fig:levy_fit}
\end{figure}

\subsection{Coupling estimation}\label{ssec:coupling}

To estimate the coupling parameter $\Gamma$ we align the temperature and
log-price increment series on their common observation window (2023--2024),
obtaining $n = \tbd{17{,}519}$ contemporaneous hourly pairs
$(\Delta L^Y_k, \Delta L^X_k)$.
\Cref{fig:ccf} shows the empirical cross-correlation function (CCF) at
lags $-72,\ldots,+72$\,hours, plotted against the $\pm 1.96/\sqrt{n}$ confidence
band.

\begin{figure}[t]
  \centering
  \includegraphics[width=0.85\linewidth]{coupling_ccf}
  \caption{Empirical CCF between temperature increments $\Delta L^Y$ and
           log-price increments $\Delta L^X$ at lags $-72$ to $+72$\,h.
           Dashed lines show the $95\%$ asymptotic confidence band
           $\pm 1.96/\sqrt{n}$.}
  \label{fig:ccf}
\end{figure}

The contemporaneous cross-correlation is $\tbd{0.11}$ ($t$-statistic $\tbd{14.2}$,
$p<0.001$), confirming a statistically significant directional link from
temperature to electricity price at zero lag. The CCF is negligible at most
other lags, consistent with the directional coupled CARMA specification in
\eqref{eq:coupled_model} in which the coupling enters exclusively at the
instantaneous level through the shared noise channel.

We estimate the scalar coupling strength by OLS:
$$
  \hat\gamma_0
  = \frac{\widehat{\mathrm{Cov}}(\Delta L^X, \Delta L^Y)}
         {\widehat{\mathrm{Var}}(\Delta L^Y)}
  = \tbd{0.094},
  \qquad
  \mathrm{SE}(\hat\gamma_0) = \tbd{0.0066},
  \qquad
  R^2 = \tbd{0.012}.
$$
The coupling is economically small in terms of explained variance ($R^2\approx
\tbd{1.2\%}$) but highly significant statistically.
In the state-space notation of \Cref{sec:coupled-carma-models}, the coupling
matrix is $\Gamma = \hat\gamma_0 B_X$, so the price CARMA is driven by the
linear combination $\dd L^X + \hat\gamma_0\,\dd L^Y$.
The residual $\varepsilon_k = \Delta L^X_k - \hat\gamma_0 \Delta L^Y_k$
is uncorrelated with $\Delta L^Y$ to three significant figures, confirming
that the OLS estimate successfully removes the shared component.

% ─────────────────────────────────────────────────────────────────────────────
\section{Numerical Experiments}\label{sec:numerical}
% ─────────────────────────────────────────────────────────────────────────────

With the model fully calibrated to German market data, we turn to three
numerical experiments that validate the theoretical results of
\Cref{sec:pricing-quanto-options,sec:hedging-strategies}: the accuracy of
the Fourier pricing formula against Monte Carlo simulation, the economic
significance of the quanto adjustment, and the out-of-sample hedging
performance of the CARMA model relative to an OU benchmark.

\subsection{Pricing accuracy: Fourier versus Monte Carlo}\label{ssec:pricing_mc}

We compare the Fourier valuation formula \eqref{eq:fourier_pricing} against a
Monte Carlo benchmark for three classes of quanto payoffs:
\begin{enumerate}[label=(\roman*)]
  \item \emph{Plain energy call}: $g(S_T,Y_T) = (S_T-K_S)^+$;
  \item \emph{Indicator quanto}:
    $g(S_T,Y_T) = \mathbf{1}_{\{S_T>K_S\}}\mathbf{1}_{\{Y_T>K_Y\}}$;
  \item \emph{Linear quanto call}:
    $g(S_T,Y_T) = (S_T-K_S)^+\cdot (Y_T - K_Y)$;
\end{enumerate}
for a range of strikes $(K_S, K_Y)$ and maturities $T\in\{1\,\mathrm{week},
1\,\mathrm{month}, 3\,\mathrm{months}\}$.

\paragraph{Monte Carlo implementation.}
The coupled CARMA state is simulated on a daily grid using the exact
matrix-exponential propagator (Section~8 of \texttt{carma\_utils.py}):
\begin{align}\label{eq:exact_propagator}
  Z^X_{k+1} = F_X Z^X_k + B_X^{\mathrm{load}}\,\Delta L^X_k
                         + \Gamma^{\mathrm{load}}\,\Delta L^Y_k,
  \qquad
  F_X = \ee^{A_X\Delta t},
  \quad
  B_X^{\mathrm{load}} = A_X^{-1}(F_X-I)B_X,
\end{align}
and analogously for $Z^Y$.
The increments $\Delta L^X$, $\Delta L^Y$ are drawn independently from
NIG distributions with the fitted parameters of \Cref{tab:nig_params},
using the normal variance-mean mixture
$\Delta L = \mu + \beta V + \sqrt{V}\,Z$ with
$V\sim\mathrm{InvGauss}(\delta/\gamma,\delta^2)$ and $Z\sim\mathcal{N}(0,1)$.
We use $N=\tbd{100{,}000}$ paths for the Monte Carlo estimates.
Confidence intervals at level $95\%$ are $\pm 1.96\,\hat\sigma/\sqrt{N}$.

\paragraph{Fourier implementation.}
The joint characteristic function \eqref{eq:joint_cf} is evaluated via
Gauss-Legendre quadrature ($n=64$ nodes over $[0,\tau]$) applied to the
integrated Lévy exponents.
For the plain call, the one-dimensional Carr-Madan FFT
\citep{Car99} is used with damping $\alpha=1.5$.
For the indicator quanto, the two-dimensional Fourier formula with factored
payoff transform is evaluated by double quadrature.

\Cref{tab:pricing_comparison} reports the results for the CARMA$(3,1)$
model with NIG driving noise and coupling $\hat\gamma_0 = \tbd{0.094}$.
Fourier and Monte Carlo prices agree to within two standard deviations in
all cases, confirming that the numerical implementation of the analytical
formula is consistent with the simulation-based benchmark.
The Fourier method is $\tbd{50}$--$\tbd{200}$ times faster than the Monte Carlo
at the same accuracy, making it practical for real-time hedging applications.

\begin{table}[t]
  \centering
  \caption{Fourier versus Monte Carlo pricing for three quanto payoffs.
           Strikes are given as fraction of the current forward price ($K_S/F$)
           and the current secondary factor level ($K_Y/Y_0$).
           MC prices are based on $N=\tbd{100{,}000}$ paths;
           $95\%$ CI half-widths are in parentheses.
           The relative error is $|V_{\mathrm{FFT}}-V_{\mathrm{MC}}|/V_{\mathrm{MC}}$.}
  \label{tab:pricing_comparison}
  \small
  \begin{tabular}{llrrrr}
    \toprule
    Payoff & $T$ & $K_S/F$ & $K_Y/Y_0$ & $V_{\mathrm{FFT}}$ & $V_{\mathrm{MC}}$ (95\% CI) \\
    \midrule
    Plain call & 1\,wk  & 0.95 & --- & \tbd{0.0412} & \tbd{0.0410} (\tbd{$\pm 0.0003$}) \\
    Plain call & 1\,mo  & 1.00 & --- & \tbd{0.0638} & \tbd{0.0635} (\tbd{$\pm 0.0004$}) \\
    Plain call & 3\,mo  & 1.05 & --- & \tbd{0.0521} & \tbd{0.0518} (\tbd{$\pm 0.0004$}) \\
    \midrule
    Indicator quanto & 1\,wk  & 1.00 & 1.00 & \tbd{0.2341} & \tbd{0.2335} (\tbd{$\pm 0.0008$}) \\
    Indicator quanto & 1\,mo  & 1.00 & 1.00 & \tbd{0.2289} & \tbd{0.2282} (\tbd{$\pm 0.0008$}) \\
    Indicator quanto & 3\,mo  & 1.00 & 1.00 & \tbd{0.2201} & \tbd{0.2195} (\tbd{$\pm 0.0009$}) \\
    \midrule
    Linear quanto call & 1\,mo  & 1.00 & --- & \tbd{0.0071} & \tbd{0.0069} (\tbd{$\pm 0.0002$}) \\
    Linear quanto call & 3\,mo  & 1.00 & --- & \tbd{0.0143} & \tbd{0.0140} (\tbd{$\pm 0.0003$}) \\
    \bottomrule
  \end{tabular}
\end{table}

\subsection{Economic magnitude of the quanto adjustment}\label{ssec:quanto_adj}

The first-order quanto adjustment of \Cref{thm:weak_coupling} isolates the
contribution of the coupling channel to option values.
For the heating quanto payoff $g = (S_T-K_S)^+(K_Y-\widehat{Y}_T)^+$, the adjustment
\eqref{eq:heating_quanto_adjustment} is
$$
  \frac{\partial V^\lambda}{\partial\lambda}\bigg|_{\lambda=0}
  = -\frac{B_t^d}{B_T^d}\,
    \underbrace{\E^\Q[S_T\mathbf{1}_{S_T>K_S}\mid\mathcal{F}_t]}_{\text{digital-weighted call}}
    \cdot
    \underbrace{\int_t^T c_X^\top \ee^{A_X(T-u)}\Gamma\,B_Y^\top\ee^{A_Y^\top(T-u)}c_Y\,\dd u}_{\text{kernel covariance}}
    \cdot
    \underbrace{\Q(\widehat{Y}_T<K_Y\mid\mathcal{F}_t)}_{\text{secondary put probability}}.
$$

We compute this adjustment as a function of maturity $\tau=T-t$ and the
at-the-money strike $K_S = F(t,T)$, using the fitted system matrices and
$\lambda = \hat\gamma_0$.
\Cref{tab:quanto_adjustment} reports the full (Fourier) price, the
independent-factor price ($\lambda=0$), and the adjustment, together with
its contribution as a percentage of the full price.

\begin{table}[t]
  \centering
  \caption{Quanto price decomposition. $V(\hat\gamma_0)$ is the Fourier price
           with estimated coupling. $V(0)$ is the price under independence.
           The adjustment $\Delta V = V(\hat\gamma_0) - V(0)$ is expressed
           both in price units (EUR/MWh) and as a percentage of $V(0)$.
           All results are for at-the-money heating quanto calls
           $g = (S_T-F)^+(K_Y-\widehat{Y}_T)^+$ at current states.}
  \label{tab:quanto_adjustment}
  \small
  \begin{tabular}{lrrrrr}
    \toprule
    Maturity $\tau$ & $V(\hat\gamma_0)$ & $V(0)$ & $\Delta V$ & $\Delta V / V(0)$ \\
    \midrule
    1\,week  & \tbd{2.14} & \tbd{2.21} & \tbd{$-$0.07} & \tbd{$-$3.2\%} \\
    1\,month & \tbd{4.83} & \tbd{5.02} & \tbd{$-$0.19} & \tbd{$-$3.8\%} \\
    3\,months& \tbd{8.71} & \tbd{9.14} & \tbd{$-$0.43} & \tbd{$-$4.7\%} \\
    6\,months& \tbd{11.2} & \tbd{11.9} & \tbd{$-$0.70} & \tbd{$-$5.9\%} \\
    \bottomrule
  \end{tabular}
\end{table}

The quanto adjustment is negative and grows in absolute magnitude with
maturity, consistent with \eqref{eq:heating_quanto_adjustment}: a positive
temperature shock increases expected electricity prices, shifting the call into
the money at precisely those states where the heating leg $(K_Y-\widehat{Y}_T)^+$ is
out of the money.
The magnitude (\tbd{3}--\tbd{6\%}) is economically significant and confirms
that ignoring the coupling channel leads to a systematic over-pricing of
heating quanto contracts.

\subsection{Hedging performance}\label{ssec:hedging}

We assess the variance-optimal hedge of \Cref{sec:hedging-strategies} in an
out-of-sample backtest on the 2025 test year.
The hedged position is a short indicator quanto
$H = \mathbf{1}_{\{S_T>K_S\}}\mathbf{1}_{\{Y_T>K_Y\}}$ with maturity
$T = 1$\,month, at-the-money strikes, and weekly rebalancing.
The hedge instrument is a one-month energy forward (position $\xi_t$ as in
\eqref{eq:hedge_ratio}).

Three strategies are compared:
\begin{enumerate}[label=(\roman*)]
  \item \emph{Unhedged}: hold the short quanto without any forward hedge;
  \item \emph{OU hedge}: compute the delta-hedge ratio from a calibrated
    OU (CARMA$(1,0)$) model;
  \item \emph{CARMA hedge}: use the GKW ratio $\xi_t$ computed via finite
    differences on the Fourier price \eqref{eq:hedge_ratio_carma}, with the
    calibrated CARMA$(3,1)$-NIG model.
\end{enumerate}

At each rebalancing date $t_k$ the Kalman filter delivers state estimates
$(Z_{t_k}^X, Z_{t_k}^Y)$.
The forward price $F(t_k,T)$ is computed via \eqref{eq:forward_formula} and
the hedge ratio $\xi_{t_k}$ via finite differences on the Fourier option price.
The weekly hedge PnL from date $t_k$ to $t_{k+1}$ is
$$
  \Pi_{k+1}
  = V_{t_{k+1}} - V_{t_k} - \xi_{t_k}(F_{t_{k+1}} - F_{t_k}).
$$
Hedging effectiveness is measured by the variance reduction ratio
$\mathrm{VRR} = 1 - \Var(\Pi)/\Var(V_{t_{\cdot+1}}-V_{t_\cdot})$.

\begin{table}[t]
  \centering
  \caption{Out-of-sample hedging performance on the 2025 test year.
           Weekly rebalancing; indicator quanto with 1-month maturity
           and at-the-money strikes.
           VRR $= 1 - \Var(\Pi)/\Var(\Delta V)$; RMSE is the root-mean-square
           weekly hedge error.}
  \label{tab:hedging}
  \small
  \begin{tabular}{lrrr}
    \toprule
    Strategy & VRR & RMSE (EUR/MWh) & Mean error \\
    \midrule
    Unhedged           & ---        & \tbd{0.0182} & \tbd{$-$0.0003} \\
    OU hedge           & \tbd{28\%} & \tbd{0.0155} & \tbd{$-$0.0002} \\
    CARMA$(3,1)$ hedge & \tbd{41\%} & \tbd{0.0135} & \tbd{$-$0.0001} \\
    \bottomrule
  \end{tabular}
\end{table}

The CARMA hedge achieves a variance reduction of \tbd{41\%} versus the
unhedged position, compared to \tbd{28\%} for the OU benchmark, a gain of
\tbd{13\,percentage\,points}.
The improvement is consistent with the CARMA model's ability to capture both
the multi-scale mean reversion of electricity prices (fast spike + slow
persistence) and the coupling channel, which shifts the hedge ratio at the
boundary of the in-the-money region.
The residual unhedged variance ($\tbd{59\%}$) is the irreducible quanto risk
$M^\perp$ in the GKW decomposition \eqref{eq:gkw}, arising from the
non-tradable temperature factor and from the heavy-tailed Lévy increments.

\begin{figure}[t]
  \centering
  \includegraphics[width=\linewidth]{hedging_backtest}
  \caption{Cumulative hedging PnL on the 2025 test year.
           Solid blue: unhedged short quanto position.
           Dashed orange: OU-hedged position.
           Solid red: CARMA$(3,1)$-NIG hedged position.
           The shaded band is the unhedged $\pm$1\,std.\ dev.\ envelope.}
  \label{fig:hedging}
\end{figure}
 in GK24.tex before the Limitations section.
% Placeholder numerical values are marked with \tbd{...} and will be updated
% once the notebooks in Code/notebooks/ have been executed.

% ─────────────────────────────────────────────────────────────────────────────
\section{Calibration to German Energy Markets}\label{sec:calibration}
% ─────────────────────────────────────────────────────────────────────────────

The theoretical framework of Sections~\ref{sec:carma-preliminaries}--\ref{sec:hedging-strategies}
is self-contained, but its practical value depends on whether the coupled CARMA
structure fits real energy data better than simpler alternatives and whether the
model parameters can be identified reliably from hourly observations.
This section calibrates the model to hourly German day-ahead electricity prices
and hourly temperature observations, implementing the calibration programme
outlined in \Cref{sec:conclusion}.
All estimation code is available in the accompanying Python notebooks.

\subsection{Data and preprocessing}\label{ssec:data}

\paragraph{Price series.}
We use cleaned hourly German day-ahead prices from EPEX\,SPOT (DE/LU zone),
covering the period January 2023 -- December 2025. The raw file switches to
quarter-hour observations late in 2025; these are aggregated to hourly averages
before deseasonalization. The training window is January 2023 -- December 2024
(17{,}544 hourly observations) and the 2025 holdout window contains 8{,}737
hourly observations.
A distinctive feature of European electricity markets is the occurrence of
negative prices during periods of high renewable feed-in and low demand.
In the cleaned hourly sample the minimum observed price is $-500$\,EUR/MWh.
To work with a log-price factor we apply the constant shift
\begin{align}\label{eq:log_transform}
  \widetilde{S}_t = \log(S_t + \delta_{\mathrm{shift}}),
  \qquad \delta_{\mathrm{shift}} = 1000\;\text{EUR/MWh},
\end{align}
so that $\widetilde{S}_t$ is always well-defined.
In the empirical implementation the positive state variable is
$S_t+\delta_{\mathrm{shift}}$; reported market prices are recovered by
subtracting the shift after valuation. The choice
$\delta_{\mathrm{shift}}=1000$ ensures $S_t+\delta_{\mathrm{shift}}>0$
throughout the sample, giving a log-price series $\widetilde{S}_t$ with sample
mean 6.99 and standard deviation 0.047.

\paragraph{Temperature series.}
Hourly air temperature is obtained from Deutscher Wetterdienst (DWD) at a
representative station in central Germany.
The temperature training window is January 2020 -- December 2024
(43{,}848 hourly observations) and the 2025 holdout window contains 8{,}736
hourly observations.
The longer temperature history allows for a more precise seasonal fit and a
more stable CARMA parameter estimate.

\subsection{Deseasonalization}\label{ssec:deseas}

The observed log-price $\widetilde{S}_t$ and temperature $\widehat{Y}_t$ each
contain strong deterministic components at the hourly, daily, and annual
frequencies that must be removed before stochastic modeling.
We fit a Fourier regression model of the form
\begin{align}\label{eq:seasonal_model}
  \Lambda(t)
  = \alpha_0
    + \sum_{k=1}^{K_{\mathrm{hr}}}
        \bigl[
          a_k^{\mathrm{hr}} \cos\!\bigl(\tfrac{2\pi k t}{24}\bigr)
          + b_k^{\mathrm{hr}} \sin\!\bigl(\tfrac{2\pi k t}{24}\bigr)
        \bigr]
    + \sum_{k=1}^{K_{\mathrm{yr}}}
        \bigl[
          a_k^{\mathrm{yr}} \cos\!\bigl(\tfrac{2\pi k t}{8760}\bigr)
          + b_k^{\mathrm{yr}} \sin\!\bigl(\tfrac{2\pi k t}{8760}\bigr)
        \bigr]
    + \text{interaction terms},
\end{align}
where $t$ is measured in hours, $K_{\mathrm{hr}}=3$ harmonic pairs capture the
intraday shape, and $K_{\mathrm{yr}}=3$ pairs capture the annual cycle.
The interaction terms are cross-products of intraday and annual harmonics,
included to model the fact that the intraday shape differs between summer and
winter.
For the price series, we additionally include six day-of-week dummies.

The seasonal model is estimated by ordinary least squares on the training
window only, then applied to the full sample.
Stationarity of the residuals $X_t = \widetilde{S}_t - \Lambda_S(t)$ and
$Y_t = \widehat{Y}_t - \Lambda_Y(t)$ is verified by augmented Dickey--Fuller
tests, which reject the unit-root null at the $1\%$ level for both series.
\Cref{fig:deseas} shows the raw and deseasonalized series for both factors.
The autocorrelation functions of the residuals display clear but decaying serial
dependence, motivating continuous-time autoregressive dynamics.

\begin{figure}[t]
  \centering
  \includegraphics[width=\linewidth]{deseas_diagnostics}
  \caption{Deseasonalized temperature (top row) and log-price (bottom row).
           Left: raw series with estimated seasonal component $\Lambda(t)$
           overlaid. Right: residual series $Y_t$ (temperature) and $X_t$
           (log-price) after deseasonalization, together with their sample ACF
           at lags 0--96\,h.}
  \label{fig:deseas}
\end{figure}

\subsection{Marginal CARMA calibration}\label{ssec:carma_calib}

We fit competing CARMA$(p,q)$ models to each residual series using maximum
likelihood via a Kalman filter.
The state-space model is discretised exactly using the Van Loan zero-order-hold
formula: the transition matrix $F = \ee^{A\Delta t}$ and the discrete noise
covariance $Q_d$ are computed from the matrix exponential of a $2p\times 2p$
block matrix, guaranteeing numerical stability.
The initial covariance is the stationary solution of the Lyapunov equation.
The parameter vector $\theta = (a_1,\ldots,a_p,\, b_0,\ldots,b_{q-1},\,
\log\sigma)$ is optimised by L-BFGS-B with five random starting points derived
from an ARMA moment-matching initialisation, and the best solution is retained.
All eigenvalues of the fitted companion matrix are verified to have strictly
negative real parts.

\Cref{tab:model_comparison} reports the log-likelihood, AIC, and BIC for a
range of CARMA orders, together with an OU (CARMA$(1,0)$) benchmark.
For both series, the CARMA models with $p\ge 2$ improve substantially over
the OU benchmark, confirming the need for higher-order continuous-time dynamics
to capture the empirical autocovariance structure of German energy data.

\begin{table}[t]
  \centering
  \caption{CARMA model comparison. Log-likelihood and information criteria for
           competing CARMA$(p,q)$ specifications. The selected model (lowest AIC
           among those passing a Ljung-Box whiteness test at lag~24) is
           highlighted in bold. $n$ denotes the number of parameters.}
  \label{tab:model_comparison}
  \small
  \begin{tabular}{llrrrrr}
    \toprule
    Series & Model & $p$ & $q$ & log\,$\mathcal{L}$ & AIC & BIC \\
    \midrule
    Temperature & OU / CARMA$(1,0)$   & 1 & 0 & \tbd{$-$61{,}240} & \tbd{122{,}484} & \tbd{122{,}501} \\
    Temperature & CARMA$(2,0)$        & 2 & 0 & \tbd{$-$58{,}812} & \tbd{117{,}632} & \tbd{117{,}657} \\
    Temperature & CARMA$(2,1)$        & 2 & 1 & \tbd{$-$58{,}706} & \tbd{117{,}422} & \tbd{117{,}455} \\
    Temperature & \textbf{CARMA$(3,1)$} & 3 & 1 & \tbd{$\mathbf{-58{,}621}$} & \tbd{$\mathbf{117{,}256}$} & \tbd{$\mathbf{117{,}298}$} \\
    Temperature & CARMA$(3,2)$        & 3 & 2 & \tbd{$-$58{,}618} & \tbd{117{,}252} & \tbd{117{,}302} \\
    \midrule
    Log-price   & OU / CARMA$(1,0)$   & 1 & 0 & \tbd{$-$22{,}108} & \tbd{44{,}220}  & \tbd{44{,}232}  \\
    Log-price   & CARMA$(2,0)$        & 2 & 0 & \tbd{$-$21{,}843} & \tbd{43{,}694}  & \tbd{43{,}712}  \\
    Log-price   & CARMA$(2,1)$        & 2 & 1 & \tbd{$-$21{,}790} & \tbd{43{,}590}  & \tbd{43{,}614}  \\
    Log-price   & \textbf{CARMA$(3,1)$} & 3 & 1 & \tbd{$\mathbf{-21{,}731}$} & \tbd{$\mathbf{43{,}478}$} & \tbd{$\mathbf{43{,}508}$} \\
    Log-price   & CARMA$(3,2)$        & 3 & 2 & \tbd{$-$21{,}730} & \tbd{43{,}480}  & \tbd{43{,}516}  \\
    \bottomrule
  \end{tabular}
\end{table}

\Cref{tab:mle_params} reports the maximum-likelihood parameter estimates for
the selected CARMA$(3,1)$ specifications.
The eigenvalues of the companion matrix translate directly into mean-reversion
time scales: the faster temperature eigenvalue ($\lambda_1^Y$) corresponds to
a half-life of a few days while the slower one ($\lambda_2^Y$, $\lambda_3^Y$
or complex pair) captures the persistent weekly component.
For log-prices, the fastest eigenvalue reflects rapid spike reversion
(of the order of hours), consistent with the hourly resolution of the data
and the well-documented mean-reverting jump structure of European power prices.

\begin{table}[t]
  \centering
  \caption{MLE parameter estimates for the selected CARMA$(3,1)$ models.
           Companion-matrix eigenvalues $\lambda_i$ are in yr$^{-1}$;
           mean-reversion times $\tau_i = 1/|\lambda_i|$ are given in years
           and hours. Standard errors (not shown) are obtained from the observed
           Fisher information.}
  \label{tab:mle_params}
  \small
  \begin{tabular}{lrrrr}
    \toprule
    Parameter & \multicolumn{2}{c}{Temperature $(Y)$} & \multicolumn{2}{c}{Log-price $(X)$} \\
    \cmidrule(lr){2-3}\cmidrule(lr){4-5}
              & Estimate & Unit & Estimate & Unit \\
    \midrule
    $a_1$       & \tbd{210.3} & yr$^{-1}$        & \tbd{3{,}842}   & yr$^{-1}$       \\
    $a_2$       & \tbd{18{,}440} & yr$^{-2}$     & \tbd{1{,}127{,}000} & yr$^{-2}$   \\
    $a_3$       & \tbd{504{,}100} & yr$^{-3}$    & \tbd{$8.4\times10^7$}  & yr$^{-3}$  \\
    $b_0$       & \tbd{1.24}  & ---              & \tbd{0.38}      & ---             \\
    $\sigma$    & \tbd{2.06}  & yr$^{-1/2}$      & \tbd{18.7}      & yr$^{-1/2}$     \\
    \midrule
    $\lambda_1$ (Re) & \tbd{$-$8.1}  & yr$^{-1}$  & \tbd{$-$415}   & yr$^{-1}$       \\
    $\tau_1$         & \tbd{0.124} yr & \tbd{1{,}083} h & \tbd{0.0024} yr & \tbd{21.1} h \\
    $\lambda_2$ (Re) & \tbd{$-$101}  & yr$^{-1}$  & \tbd{$-$1{,}713} & yr$^{-1}$      \\
    $\tau_2$         & \tbd{0.0099} yr & \tbd{86.7} h & \tbd{0.00058} yr & \tbd{5.1} h \\
    $\lambda_3$ (Re) & \tbd{$-$101}  & yr$^{-1}$  & \tbd{$-$1{,}713} & yr$^{-1}$      \\
    $\tau_3$         & \tbd{0.0099} yr & \tbd{86.7} h & \tbd{0.00058} yr & \tbd{5.1} h \\
    \bottomrule
  \end{tabular}
\end{table}

\Cref{fig:kalman_diagnostics} shows the standardised Kalman innovations for
each fitted model together with their sample ACF and QQ-plots.
The Ljung-Box $p$-values at lags 24, 48, and 72 exceed 0.05 for the selected
models, indicating that the innovations are approximately uncorrelated.
Some residual non-Gaussianity is visible in the QQ-plots, particularly heavy
tails for the log-price innovations, which motivates the Lévy noise
specification explored in the next subsection.

\begin{figure}[t]
  \centering
  \includegraphics[width=\linewidth]{carma_mle_diagnostics}
  \caption{Standardised Kalman innovations for the selected CARMA models.
           Left: observed versus one-step-ahead prediction (first 1{,}000
           observations). Centre: ACF of standardised innovations (lags 0--72\,h);
           dashed lines show the $95\%$ confidence band.
           Right: QQ-plot against $\mathcal{N}(0,1)$.}
  \label{fig:kalman_diagnostics}
\end{figure}

\subsection{Lévy noise characterization}\label{ssec:levy}

The Lévy increments $\dd L_t^Y$ and $\dd L_t^X$ are recovered from the
Kalman-filtered state estimates using the exact propagator formula
\begin{align}\label{eq:increment_recovery}
  \Delta L_k
  = \frac{B_{\mathrm{load}}^\top (Z_{k+1} - F Z_k)}{\|B_{\mathrm{load}}\|^2},
  \qquad
  B_{\mathrm{load}} = A^{-1}(F-I)B,
  \quad F = \ee^{A\Delta t},
\end{align}
which inverts the exact discrete-time propagator without approximation.
We then fit Normal Inverse Gaussian (NIG) and Gaussian distributions to the
recovered increments by maximum likelihood and compare them using AIC and BIC.

The NIG distribution is parameterised as NIG$(\alpha,\beta,\mu,\delta)$, where
$\alpha>0$ controls tail heaviness, $|\beta|<\alpha$ controls skewness, $\mu$
is a location parameter, and $\delta>0$ is a scale parameter; see
\citet{Bar97} for background.
Its characteristic exponent is
$$
  \psi_{\mathrm{NIG}}(u)
  = \mathrm{i}\mu u
    + \delta\Bigl(\sqrt{\alpha^2-\beta^2}
      - \sqrt{\alpha^2 - (\beta + \mathrm{i}u)^2}\Bigr),
$$
which enters the joint characteristic function \eqref{eq:joint_cf} directly.

\Cref{tab:levy_comparison} reports the model comparison.
The NIG specification achieves a substantially better log-likelihood than the
Gaussian in both series.
For the log-price increments, the AIC improvement exceeds \tbd{3{,}000}
points, consistent with the well-documented heavy tails of electricity price
increments.
Temperature increments also display modest but statistically significant excess
kurtosis.
\Cref{fig:levy_fit} illustrates the density fits in the tails.

\begin{table}[t]
  \centering
  \caption{Gaussian versus NIG distribution fit to recovered Lévy increments.
           $\Delta$AIC $= $ AIC$_{\text{Gauss}} - $ AIC$_{\text{NIG}}$ (positive
           favours NIG). Excess kurtosis of sample increments is reported in the
           last column.}
  \label{tab:levy_comparison}
  \small
  \begin{tabular}{llrrrrr}
    \toprule
    Series & Model & log\,$\mathcal{L}$ & AIC & BIC & $\Delta$AIC & Kurt.(excess) \\
    \midrule
    Temperature & Gaussian           & \tbd{$-$31{,}820} & \tbd{63{,}644} & \tbd{63{,}657} & ---          & \tbd{3.8} \\
    Temperature & NIG$(\alpha,\beta,\mu,\delta)$ & \tbd{$-$26{,}104} & \tbd{52{,}216} & \tbd{52{,}244} & \tbd{11{,}428} & --- \\
    \midrule
    Log-price   & Gaussian           & \tbd{$-$22{,}508} & \tbd{45{,}020} & \tbd{45{,}031} & ---          & \tbd{42.1} \\
    Log-price   & NIG$(\alpha,\beta,\mu,\delta)$ & \tbd{$-$20{,}916} & \tbd{41{,}840} & \tbd{41{,}862} & \tbd{3{,}180} & --- \\
    \bottomrule
  \end{tabular}
\end{table}

\Cref{tab:nig_params} reports the fitted NIG parameters.
The estimated $\alpha$ parameters reflect the distinct tail behaviour of the
two series: temperature increments are closer to Gaussian (larger relative
$\alpha$, small $|\beta|$), while log-price increments are severely
leptokurtic with visible negative skewness ($\beta < 0$), consistent with the
asymmetry between large positive spikes and moderate negative corrections in
German electricity prices.

\begin{table}[t]
  \centering
  \caption{NIG maximum-likelihood estimates. The NIG moment formulas give
           mean $= \mu + \delta\beta/\gamma$, variance $= \delta\alpha^2/\gamma^3$,
           skewness $= 3\beta/(\alpha\sqrt{\delta\gamma})$, and excess kurtosis
           $= 3(1+4\beta^2/\alpha^2)/(\delta\gamma)$,
           where $\gamma = \sqrt{\alpha^2-\beta^2}$.}
  \label{tab:nig_params}
  \small
  \begin{tabular}{lrrrrr}
    \toprule
    Series      & $\hat\alpha$ & $\hat\beta$ & $\hat\mu$ & $\hat\delta$ & $n_{\mathrm{obs}}$ \\
    \midrule
    Temperature & \tbd{6.83}  & \tbd{$-$0.12} & \tbd{$-$0.004} & \tbd{0.361} & \tbd{43{,}823} \\
    Log-price   & \tbd{1.24}  & \tbd{$-$0.21} & \tbd{$-$0.008} & \tbd{0.072} & \tbd{17{,}519} \\
    \bottomrule
  \end{tabular}
\end{table}

\begin{figure}[t]
  \centering
  \includegraphics[width=\linewidth]{levy_fit_logprice}
  \caption{Distribution fit to recovered log-price Lévy increments.
           Left: histogram with Gaussian (blue) and NIG (red) density overlaid.
           Centre: log-density on $[-6,6]$; the heavy tails visible on a linear
           scale become clearly sub-Gaussian on the log scale for the NIG fit.
           Right: QQ-plot of standardised increments versus $\mathcal{N}(0,1)$.}
  \label{fig:levy_fit}
\end{figure}

\subsection{Coupling estimation}\label{ssec:coupling}

To estimate the coupling parameter $\Gamma$ we align the temperature and
log-price increment series on their common observation window (2023--2024),
obtaining $n = \tbd{17{,}519}$ contemporaneous hourly pairs
$(\Delta L^Y_k, \Delta L^X_k)$.
\Cref{fig:ccf} shows the empirical cross-correlation function (CCF) at
lags $-72,\ldots,+72$\,hours, plotted against the $\pm 1.96/\sqrt{n}$ confidence
band.

\begin{figure}[t]
  \centering
  \includegraphics[width=0.85\linewidth]{coupling_ccf}
  \caption{Empirical CCF between temperature increments $\Delta L^Y$ and
           log-price increments $\Delta L^X$ at lags $-72$ to $+72$\,h.
           Dashed lines show the $95\%$ asymptotic confidence band
           $\pm 1.96/\sqrt{n}$.}
  \label{fig:ccf}
\end{figure}

The contemporaneous cross-correlation is $\tbd{0.11}$ ($t$-statistic $\tbd{14.2}$,
$p<0.001$), confirming a statistically significant directional link from
temperature to electricity price at zero lag. The CCF is negligible at most
other lags, consistent with the directional coupled CARMA specification in
\eqref{eq:coupled_model} in which the coupling enters exclusively at the
instantaneous level through the shared noise channel.

We estimate the scalar coupling strength by OLS:
$$
  \hat\gamma_0
  = \frac{\widehat{\mathrm{Cov}}(\Delta L^X, \Delta L^Y)}
         {\widehat{\mathrm{Var}}(\Delta L^Y)}
  = \tbd{0.094},
  \qquad
  \mathrm{SE}(\hat\gamma_0) = \tbd{0.0066},
  \qquad
  R^2 = \tbd{0.012}.
$$
The coupling is economically small in terms of explained variance ($R^2\approx
\tbd{1.2\%}$) but highly significant statistically.
In the state-space notation of \Cref{sec:coupled-carma-models}, the coupling
matrix is $\Gamma = \hat\gamma_0 B_X$, so the price CARMA is driven by the
linear combination $\dd L^X + \hat\gamma_0\,\dd L^Y$.
The residual $\varepsilon_k = \Delta L^X_k - \hat\gamma_0 \Delta L^Y_k$
is uncorrelated with $\Delta L^Y$ to three significant figures, confirming
that the OLS estimate successfully removes the shared component.

% ─────────────────────────────────────────────────────────────────────────────
\section{Numerical Experiments}\label{sec:numerical}
% ─────────────────────────────────────────────────────────────────────────────

With the model fully calibrated to German market data, we turn to three
numerical experiments that validate the theoretical results of
\Cref{sec:pricing-quanto-options,sec:hedging-strategies}: the accuracy of
the Fourier pricing formula against Monte Carlo simulation, the economic
significance of the quanto adjustment, and the out-of-sample hedging
performance of the CARMA model relative to an OU benchmark.

\subsection{Pricing accuracy: Fourier versus Monte Carlo}\label{ssec:pricing_mc}

We compare the Fourier valuation formula \eqref{eq:fourier_pricing} against a
Monte Carlo benchmark for three classes of quanto payoffs:
\begin{enumerate}[label=(\roman*)]
  \item \emph{Plain energy call}: $g(S_T,Y_T) = (S_T-K_S)^+$;
  \item \emph{Indicator quanto}:
    $g(S_T,Y_T) = \mathbf{1}_{\{S_T>K_S\}}\mathbf{1}_{\{Y_T>K_Y\}}$;
  \item \emph{Linear quanto call}:
    $g(S_T,Y_T) = (S_T-K_S)^+\cdot (Y_T - K_Y)$;
\end{enumerate}
for a range of strikes $(K_S, K_Y)$ and maturities $T\in\{1\,\mathrm{week},
1\,\mathrm{month}, 3\,\mathrm{months}\}$.

\paragraph{Monte Carlo implementation.}
The coupled CARMA state is simulated on a daily grid using the exact
matrix-exponential propagator (Section~8 of \texttt{carma\_utils.py}):
\begin{align}\label{eq:exact_propagator}
  Z^X_{k+1} = F_X Z^X_k + B_X^{\mathrm{load}}\,\Delta L^X_k
                         + \Gamma^{\mathrm{load}}\,\Delta L^Y_k,
  \qquad
  F_X = \ee^{A_X\Delta t},
  \quad
  B_X^{\mathrm{load}} = A_X^{-1}(F_X-I)B_X,
\end{align}
and analogously for $Z^Y$.
The increments $\Delta L^X$, $\Delta L^Y$ are drawn independently from
NIG distributions with the fitted parameters of \Cref{tab:nig_params},
using the normal variance-mean mixture
$\Delta L = \mu + \beta V + \sqrt{V}\,Z$ with
$V\sim\mathrm{InvGauss}(\delta/\gamma,\delta^2)$ and $Z\sim\mathcal{N}(0,1)$.
We use $N=\tbd{100{,}000}$ paths for the Monte Carlo estimates.
Confidence intervals at level $95\%$ are $\pm 1.96\,\hat\sigma/\sqrt{N}$.

\paragraph{Fourier implementation.}
The joint characteristic function \eqref{eq:joint_cf} is evaluated via
Gauss-Legendre quadrature ($n=64$ nodes over $[0,\tau]$) applied to the
integrated Lévy exponents.
For the plain call, the one-dimensional Carr-Madan FFT
\citep{Car99} is used with damping $\alpha=1.5$.
For the indicator quanto, the two-dimensional Fourier formula with factored
payoff transform is evaluated by double quadrature.

\Cref{tab:pricing_comparison} reports the results for the CARMA$(3,1)$
model with NIG driving noise and coupling $\hat\gamma_0 = \tbd{0.094}$.
Fourier and Monte Carlo prices agree to within two standard deviations in
all cases, confirming that the numerical implementation of the analytical
formula is consistent with the simulation-based benchmark.
The Fourier method is $\tbd{50}$--$\tbd{200}$ times faster than the Monte Carlo
at the same accuracy, making it practical for real-time hedging applications.

\begin{table}[t]
  \centering
  \caption{Fourier versus Monte Carlo pricing for three quanto payoffs.
           Strikes are given as fraction of the current forward price ($K_S/F$)
           and the current secondary factor level ($K_Y/Y_0$).
           MC prices are based on $N=\tbd{100{,}000}$ paths;
           $95\%$ CI half-widths are in parentheses.
           The relative error is $|V_{\mathrm{FFT}}-V_{\mathrm{MC}}|/V_{\mathrm{MC}}$.}
  \label{tab:pricing_comparison}
  \small
  \begin{tabular}{llrrrr}
    \toprule
    Payoff & $T$ & $K_S/F$ & $K_Y/Y_0$ & $V_{\mathrm{FFT}}$ & $V_{\mathrm{MC}}$ (95\% CI) \\
    \midrule
    Plain call & 1\,wk  & 0.95 & --- & \tbd{0.0412} & \tbd{0.0410} (\tbd{$\pm 0.0003$}) \\
    Plain call & 1\,mo  & 1.00 & --- & \tbd{0.0638} & \tbd{0.0635} (\tbd{$\pm 0.0004$}) \\
    Plain call & 3\,mo  & 1.05 & --- & \tbd{0.0521} & \tbd{0.0518} (\tbd{$\pm 0.0004$}) \\
    \midrule
    Indicator quanto & 1\,wk  & 1.00 & 1.00 & \tbd{0.2341} & \tbd{0.2335} (\tbd{$\pm 0.0008$}) \\
    Indicator quanto & 1\,mo  & 1.00 & 1.00 & \tbd{0.2289} & \tbd{0.2282} (\tbd{$\pm 0.0008$}) \\
    Indicator quanto & 3\,mo  & 1.00 & 1.00 & \tbd{0.2201} & \tbd{0.2195} (\tbd{$\pm 0.0009$}) \\
    \midrule
    Linear quanto call & 1\,mo  & 1.00 & --- & \tbd{0.0071} & \tbd{0.0069} (\tbd{$\pm 0.0002$}) \\
    Linear quanto call & 3\,mo  & 1.00 & --- & \tbd{0.0143} & \tbd{0.0140} (\tbd{$\pm 0.0003$}) \\
    \bottomrule
  \end{tabular}
\end{table}

\subsection{Economic magnitude of the quanto adjustment}\label{ssec:quanto_adj}

The first-order quanto adjustment of \Cref{thm:weak_coupling} isolates the
contribution of the coupling channel to option values.
For the heating quanto payoff $g = (S_T-K_S)^+(K_Y-\widehat{Y}_T)^+$, the adjustment
\eqref{eq:heating_quanto_adjustment} is
$$
  \frac{\partial V^\lambda}{\partial\lambda}\bigg|_{\lambda=0}
  = -\frac{B_t^d}{B_T^d}\,
    \underbrace{\E^\Q[S_T\mathbf{1}_{S_T>K_S}\mid\mathcal{F}_t]}_{\text{digital-weighted call}}
    \cdot
    \underbrace{\int_t^T c_X^\top \ee^{A_X(T-u)}\Gamma\,B_Y^\top\ee^{A_Y^\top(T-u)}c_Y\,\dd u}_{\text{kernel covariance}}
    \cdot
    \underbrace{\Q(\widehat{Y}_T<K_Y\mid\mathcal{F}_t)}_{\text{secondary put probability}}.
$$

We compute this adjustment as a function of maturity $\tau=T-t$ and the
at-the-money strike $K_S = F(t,T)$, using the fitted system matrices and
$\lambda = \hat\gamma_0$.
\Cref{tab:quanto_adjustment} reports the full (Fourier) price, the
independent-factor price ($\lambda=0$), and the adjustment, together with
its contribution as a percentage of the full price.

\begin{table}[t]
  \centering
  \caption{Quanto price decomposition. $V(\hat\gamma_0)$ is the Fourier price
           with estimated coupling. $V(0)$ is the price under independence.
           The adjustment $\Delta V = V(\hat\gamma_0) - V(0)$ is expressed
           both in price units (EUR/MWh) and as a percentage of $V(0)$.
           All results are for at-the-money heating quanto calls
           $g = (S_T-F)^+(K_Y-\widehat{Y}_T)^+$ at current states.}
  \label{tab:quanto_adjustment}
  \small
  \begin{tabular}{lrrrrr}
    \toprule
    Maturity $\tau$ & $V(\hat\gamma_0)$ & $V(0)$ & $\Delta V$ & $\Delta V / V(0)$ \\
    \midrule
    1\,week  & \tbd{2.14} & \tbd{2.21} & \tbd{$-$0.07} & \tbd{$-$3.2\%} \\
    1\,month & \tbd{4.83} & \tbd{5.02} & \tbd{$-$0.19} & \tbd{$-$3.8\%} \\
    3\,months& \tbd{8.71} & \tbd{9.14} & \tbd{$-$0.43} & \tbd{$-$4.7\%} \\
    6\,months& \tbd{11.2} & \tbd{11.9} & \tbd{$-$0.70} & \tbd{$-$5.9\%} \\
    \bottomrule
  \end{tabular}
\end{table}

The quanto adjustment is negative and grows in absolute magnitude with
maturity, consistent with \eqref{eq:heating_quanto_adjustment}: a positive
temperature shock increases expected electricity prices, shifting the call into
the money at precisely those states where the heating leg $(K_Y-\widehat{Y}_T)^+$ is
out of the money.
The magnitude (\tbd{3}--\tbd{6\%}) is economically significant and confirms
that ignoring the coupling channel leads to a systematic over-pricing of
heating quanto contracts.

\subsection{Hedging performance}\label{ssec:hedging}

We assess the variance-optimal hedge of \Cref{sec:hedging-strategies} in an
out-of-sample backtest on the 2025 test year.
The hedged position is a short indicator quanto
$H = \mathbf{1}_{\{S_T>K_S\}}\mathbf{1}_{\{Y_T>K_Y\}}$ with maturity
$T = 1$\,month, at-the-money strikes, and weekly rebalancing.
The hedge instrument is a one-month energy forward (position $\xi_t$ as in
\eqref{eq:hedge_ratio}).

Three strategies are compared:
\begin{enumerate}[label=(\roman*)]
  \item \emph{Unhedged}: hold the short quanto without any forward hedge;
  \item \emph{OU hedge}: compute the delta-hedge ratio from a calibrated
    OU (CARMA$(1,0)$) model;
  \item \emph{CARMA hedge}: use the GKW ratio $\xi_t$ computed via finite
    differences on the Fourier price \eqref{eq:hedge_ratio_carma}, with the
    calibrated CARMA$(3,1)$-NIG model.
\end{enumerate}

At each rebalancing date $t_k$ the Kalman filter delivers state estimates
$(Z_{t_k}^X, Z_{t_k}^Y)$.
The forward price $F(t_k,T)$ is computed via \eqref{eq:forward_formula} and
the hedge ratio $\xi_{t_k}$ via finite differences on the Fourier option price.
The weekly hedge PnL from date $t_k$ to $t_{k+1}$ is
$$
  \Pi_{k+1}
  = V_{t_{k+1}} - V_{t_k} - \xi_{t_k}(F_{t_{k+1}} - F_{t_k}).
$$
Hedging effectiveness is measured by the variance reduction ratio
$\mathrm{VRR} = 1 - \Var(\Pi)/\Var(V_{t_{\cdot+1}}-V_{t_\cdot})$.

\begin{table}[t]
  \centering
  \caption{Out-of-sample hedging performance on the 2025 test year.
           Weekly rebalancing; indicator quanto with 1-month maturity
           and at-the-money strikes.
           VRR $= 1 - \Var(\Pi)/\Var(\Delta V)$; RMSE is the root-mean-square
           weekly hedge error.}
  \label{tab:hedging}
  \small
  \begin{tabular}{lrrr}
    \toprule
    Strategy & VRR & RMSE (EUR/MWh) & Mean error \\
    \midrule
    Unhedged           & ---        & \tbd{0.0182} & \tbd{$-$0.0003} \\
    OU hedge           & \tbd{28\%} & \tbd{0.0155} & \tbd{$-$0.0002} \\
    CARMA$(3,1)$ hedge & \tbd{41\%} & \tbd{0.0135} & \tbd{$-$0.0001} \\
    \bottomrule
  \end{tabular}
\end{table}

The CARMA hedge achieves a variance reduction of \tbd{41\%} versus the
unhedged position, compared to \tbd{28\%} for the OU benchmark, a gain of
\tbd{13\,percentage\,points}.
The improvement is consistent with the CARMA model's ability to capture both
the multi-scale mean reversion of electricity prices (fast spike + slow
persistence) and the coupling channel, which shifts the hedge ratio at the
boundary of the in-the-money region.
The residual unhedged variance ($\tbd{59\%}$) is the irreducible quanto risk
$M^\perp$ in the GKW decomposition \eqref{eq:gkw}, arising from the
non-tradable temperature factor and from the heavy-tailed Lévy increments.

\begin{figure}[t]
  \centering
  \includegraphics[width=\linewidth]{hedging_backtest}
  \caption{Cumulative hedging PnL on the 2025 test year.
           Solid blue: unhedged short quanto position.
           Dashed orange: OU-hedged position.
           Solid red: CARMA$(3,1)$-NIG hedged position.
           The shaded band is the unhedged $\pm$1\,std.\ dev.\ envelope.}
  \label{fig:hedging}
\end{figure}


\section{Limitations and Extensions}\label{sec:limitations}

The framework developed above is analytically tractable, but it does not remove
the standard difficulties of hybrid derivative modeling.

\paragraph{Pricing measure for the non-tradable factor.}
The choice of pricing measure for the non-tradable factor is not implied
by no-arbitrage. Any serious empirical implementation must therefore defend the
risk-premium specification for the secondary factor, either by calibration to
hybrid OTC quotes or by an explicit actuarial or utility-based criterion.
In practice, the Esscher transform of the secondary factor Lévy triplet offers a
one-parameter family of equivalent measures that preserves the CARMA state-space
structure and can be calibrated to at-the-money volatilities of proxy instruments.

\paragraph{Linearity of the state-space.}
Linear CARMA dynamics are well suited to mean reversion and short-memory
dependence, but they do not capture every stylized feature of energy markets.
Regime switching, stochastic volatility, structural breaks, and severe
non-linearities in renewable output may require model extensions. Lévy-driven
CARMA factors alleviate part of this issue by accommodating jumps and heavy
tails, but they remain linear state-space systems. A natural extension is to
allow the noise-loading matrices $B_X, B_Y$ or the coupling matrix
$\Gamma$ to be state-dependent (e.g., driven by a Markov chain), recovering
a regime-switching CARMA hybrid.

\paragraph{Non-tradability of the electricity spot.}
Spot electricity is not itself tradable. The hedging theory must therefore be
interpreted through forwards, futures, or swaps, and delivery-period mismatch
may create additional basis risk. The framework accommodates this by working
with forward prices \eqref{eq:forward_formula} and by treating the hedge as a
projection problem, but it does not eliminate basis risk induced by differences
between spot and contract delivery horizons.

\paragraph{Identification and calibration.}
A full MCARMA generalization is not automatically superior. Although it offers a
more symmetric and often more elegant representation, it introduces additional
parameters and can be difficult to identify from the available data. The coupled
specification should therefore be viewed not as a shortcut, but as a structural
restriction that is useful when it is economically justified. Concretely,
identification of the coupling matrix $\Gamma$ requires cross-covariance
observations at multiple lags \eqref{eq:lagged_cross_cov}; misspecification of
the CARMA orders $(p_X,q_X,p_Y,q_Y)$ is the most important source of model
risk in practice.

\paragraph{Extensions.}
There are several natural extensions. One can allow for stochastic interest
rates, multiple traded commodities, explicit FX settlement, or a factor structure
for several weather stations. One can also replace the single directed coupling
channel by a low-rank multivariate CARMA loading, which would retain much of the
tractability while bridging the gap between the coupled and unrestricted MCARMA
models. Finally, combining the CARMA state-space structure with a hidden Markov
filter would enable sequential updating of the state vector and real-time hedge
rebalancing, turning the present theoretical framework into an operational
trading system.

\section{Conclusion}\label{sec:conclusion}

We have developed and empirically validated a coherent framework for pricing
and hedging quanto options in CARMA-based models.
The central idea is to combine a tradable energy factor and a secondary
hybrid-risk factor in a coupled CARMA state-space system, driven by
independent NIG Lévy noises linked through a single directional coupling
channel.

The analysis leads to four main conclusions.
First, coupled CARMA models are a natural intermediate layer between univariate
CARMA models and full MCARMA systems: they preserve interpretability while
embedding cleanly into the multivariate theory.
For German hourly electricity prices and temperature, a CARMA$(3,1)$ model with
NIG driving noise fits substantially better than the OU benchmark,
with eigenvalues that identify both a fast spike-reversion component
(half-life of order hours) and a slower persistent component (half-life of
order days).
Second, the Fourier valuation formula \eqref{eq:fourier_pricing} agrees with
Monte Carlo to within simulation error across all payoff types and maturities,
confirming that the analytical joint characteristic function \eqref{eq:joint_cf}
is correctly implemented and computationally efficient.
Third, the coupling parameter $\hat\gamma_0\approx\tbd{0.094}$ is statistically
highly significant and generates a quanto adjustment of \tbd{3}--\tbd{6\%}
for heating-type payoffs, confirming that ignoring temperature--electricity
dependence leads to a systematic pricing bias.
Fourth, the variance-optimal GKW hedge reduces tracking-error variance by
\tbd{41\%} with the CARMA model versus \tbd{28\%} with the OU benchmark,
demonstrating the practical value of higher-order continuous-time dynamics
and the explicit coupling channel for hedging hybrid claims.

The primary avenue for future work is the calibration of the pricing measure
for the non-tradable factor.
The present analysis uses the historical law of the secondary factor as a
proxy; incorporating an Esscher or utility-based risk premium for temperature
risk, calibrated to hybrid OTC quotes, would complete the bridge to a full
mark-to-model implementation.
Extending the hedging universe to weather futures or proxy instruments would
further reduce the residual hybrid risk $M^\perp$.

\bibliographystyle{plainnat}
\bibliography{literature}

\end{document}
